% concatenated from main.tex
\documentclass[11pt]{article}
\usepackage[margin=1in]{geometry}
\usepackage{amsthm}

\usepackage{graphicx} % Required for inserting images
\usepackage{amsmath,amssymb,amscd,mathtools}
\usepackage{dsfont}
\usepackage{xcolor}
\usepackage{todonotes}
\usepackage[normalem]{ulem} % Load ulem package
\usepackage{hyperref}

\newcommand{\R}{\mathbb{R}}
\newcommand{\N}{\mathbb{N}}
\newcommand{\Z}{\mathbb{Z}}
\newcommand{\Aut}{\mathrm{Aut}}
\newcommand{\p}{\varphi}
\newcommand{\im}{\mathrm{Im}}
\newcommand{\id}{\mathrm{id}}
\newcommand{\GENEO}{\mathrm{GENEO}}
\usepackage{cancel}

\newcommand{\patrizio}[1]{\textcolor{purple}{#1}}
\newcommand{\pat}[1]{\todo[author=Patrizio]{#1}}

% \newtheorem{teor}{Theorem}[section]
% \newtheorem{lem}[teor]{Lemma}
% \newtheorem{cor}[teor]{Corollary}
% \newtheorem{prop}[teor]{Proposition}
% \newtheorem{conj}[teor]{Statement}
\newtheorem{theorem}{Theorem}[section]
\newtheorem{lemma}[theorem]{Lemma}
\newtheorem{corollary}[theorem]{Corollary}
\newtheorem{proposition}[theorem]{Proposition}
\newtheorem{definition}[theorem]{Definition}
\newtheorem{rem}[theorem]{Remark}
\newtheorem{es}[theorem]{Example}


% Title. If the supplement option is on, then "Supplementary Material"
% is automatically inserted before the title.
\title{An Algebraic Representation Theorem for Linear GENEOs in Machine Learning}
\author{Francesco Conti\thanks{Inria Sophia Antipolis, DataShape team, France. Email: \texttt{francesco.conti@inria.fr}} \and 
Patrizio Frosini\thanks{University of Pisa, 56127 Pisa, Italy. Email: \texttt{patrizio.frosini@unipi.it}} \and 
Nicola Quercioli\thanks{Independent researcher. Email: \texttt{nicola.quercioli@gmail.com}}}
\date{}

\begin{document}

\maketitle

\begin{abstract}
Geometric and Topological Deep Learning enhance machine learning by exploiting geometric and topological structures in data. In this context, Group Equivariant Non-Expansive Operators (GENEOs) provide a principled framework for encoding symmetries and designing efficient, interpretable neural architectures. Originally introduced in Topological Data Analysis and later adopted in Deep Learning, GENEOs unify geometric and topological approaches and include, as a special case, the operator computing persistence diagrams. Their theoretical foundations rely on group actions, equivariance, and compactness properties of operator spaces. GENEOs are defined between data spaces endowed with transformation groups, referred to as perception pairs. 
While previous representation results characterized linear GENEOs acting on identical perception pairs, many applications require operators between heterogeneous ones. In this work, we introduce and prove a new algebraic representation theorem for linear GENEOs acting between different perception pairs. Under mild assumptions on data domains and group actions, this result provides a complete characterization of such operators and implicitly yields a computational method to generate all linear GENEOs between two perception pairs. The proof combines tools from the theory of stochastic matrices, group actions, and invariant measures. In our paper, we also show that the space of linear GENEOs is compact and convex, making easier their use for learning and optimization. 
Our new theorem enlarges the applicability of GENEOs, and allows for their composition with a broader class of current architectures. We show this with an application in image reconstruction.
\end{abstract}

\medskip
\noindent\textbf{Keywords:} Group equivariant non-expansive operator, topological data analysis, group action, permutant measure

\medskip
\noindent\textbf{MSC 2020:} 15B51, 20C35, 47B38, 55N31, 62R40, 68T07, 68T09

\section{Introduction}
Geometric Deep Learning \cite{Br17,BrBr21} and Topological Deep Learning \cite{Ri19,Ho19,Ca20,We21,Zia2024} are two closely related and highly active research areas that have recently emerged to advance Machine Learning through topological and geometric techniques. Within this framework lies the notion of Group Equivariant Non-Expansive Operators (GENEOs).

Originally introduced about ten years ago in Topological Data Analysis to handle invariance of persistent homology groups \cite{FrJa16}, GENEOs later entered Deep Learning and Artificial Intelligence as geometric–topological tools for designing efficient neural networks \cite{bergomi2019towards,Mi23}. 
The key idea is twofold: first, replacing complex networks with smaller ones constructed from a limited number of equivariant components; and second, shifting attention from the structure of data spaces to that of agent/observer spaces, understood in terms of functional operators. 
By encoding symmetries through group actions, these operators can drastically reduce the number of parameters in AI applications while yielding interpretable Machine Learning models \cite{BocchiFMPPGGLBF24,CoBoGiGiPeFr25}.

The connection between GENEOs and Geometric Deep Learning is most evident in their shared reliance on symmetries, group actions, and rich geometric structures, such as differentiable or Riemannian structures on manifolds. Notably, even the computation of the Laplacian and the multiplication by each circulant matrix $\mathrm{circ}(c_0,\ldots,c_{n-1})$ with $\sum_{i=0}^{n-1}|c_i|\le 1$ (to give just two simple examples) can be expressed as GENEOs. Equally compelling is the link with Topological Deep Learning, which builds on topological techniques—for instance, compactness theorems on spaces of operators—and on the remarkable fact that computing persistence diagrams in Topological Data Analysis can be viewed as a special case of a GENEO (see Section~\ref{msG}). Taken together, these observations underscore the role of GENEOs as a unifying framework bridging Geometric and Topological Deep Learning, thereby providing a cross-cutting tool for incorporating topological and geometric contributions into Machine Learning research.

Interest in GENEOs has steadily grown, leading to both a solid mathematical theory and practical applications across scientific domains (cf., e.g., the recent works \cite{Felipe2025,BoFeFr25,BoFrMi25,LaMiBoFrSo25,Ah26}). Once their usefulness in Machine Learning was established, a central challenge arose: designing equivariant operators with prescribed symmetries. In \cite{BoBoBrFrQu23}, a theorem was established that proves a one-to-one correspondence between certain measures, known as permutant measures, and linear GENEOs mapping a dataset—represented by real-valued functions on a discrete space—onto itself.

In many practical applications, however, the input and output data types differ, making the original theorem inapplicable. In this work, we extend the result to this broader setting by employing the generalized permutant measures introduced in \cite{FaFeFr2023}. Our main contribution is a new and more general representation theorem, valid for GENEOs acting between different perception pairs. 
This theorem enables the construction of all linear GENEOs between two given perception pairs, assuming that the data functions have finite domains and that the corresponding group actions are transitive.
Beyond its theoretical value, we demonstrate the practical significance of this result by showing how it can be effectively applied to enhance the performance of autoencoders.

Sections~\ref{Preliminaries} and \ref{Building} review the necessary definitions and preliminary results, setting the stage for Section~\ref{Representation}, where we introduce a representation theorem for linear group equivariant operators (Theorem~\ref{maint}). Building on this, Section~\ref{Linear} presents and proves our main result (Theorem~\ref{teo:linear}). Section~\ref{Compactness} establishes the compactness and convexity of the space of all linear GENEOs under our assumptions. Section~\ref{Application} then provides an application to autoencoders, demonstrating how the operators introduced by our new representation theorem can be effectively leveraged in machine learning. Finally, Section~\ref{Conclusions} concludes the paper by outlining future research directions that pave the way for further developments in this line of work.

\section{Preliminaries}
\label{Preliminaries}

In Sections~\ref{Preliminaries} and \ref{Building}, we recall the main concepts underlying this paper, including GEO and GENEO, as presented in \cite{FaFeFr2023}.

\subsection{Mathematical settings and GENEOs}\label{msG}
Assuming that $Z=\left\{ z_1, \dots, z_\ell \right\} = \left\{ z_k \right\}_{k=1}^\ell$ is a finite set with cardinality $\ell$, we can consider the space $\R^Z:= \{\varphi \colon Z \to \R \}$ of all real-valued functions on $Z$, equipped with the usual uniform norm $\left\| \cdot \right\|_\infty$. We recall that, if $Z$ is finite, $\R^Z \cong \R^\ell$, and $\R^Z$ has the canonical basis $\mathcal{B}_{Z}=\left( \mathds{1}_{z_1}, \dots, \mathds{1}_{z_\ell} \right)$, where for any $k = 1, \dots, \ell$ the function $\mathds{1}_{z_k} \colon X \to \R$ takes value $1$ on $z_k$, and $0$ otherwise. Moreover, $\Aut(Z)$ denotes the group of bijection from $Z$ to itself with respect to the usual composition operation. $\Aut(Z)$  acts on $\R^Z$ by composition on the right. Considering a group $G \subseteq \Aut(Z)$ we say that the couple $\left( \R^Z, G \right)$ is a perception pair.

{\color{black} To avoid overly heavy notation, in what follows we will juxtapose functions both to denote function composition and to denote their product. The context will make the meaning of the symbols clear.}

\begin{definition}
    Given two finite sets $X,Y$, we consider two perception pairs $\left(\R^X, G\right)$, $\left(\R^Y, K\right)$, where $G \subseteq \Aut(X)$ and $K\subseteq\Aut(Y)$. Each map $\left(F, T \right) \colon \left(\R^X, G\right) \to \left(\R^Y, K\right)$ such that $T$ is a homomorphism and $F$ is $T$-equivariant (i.e., $F(\p g) = F(\p)T (g)$ for every $\p \in \Phi, g \in G$) is said to be a \textbf{Group Equivariant Operator} (GEO) from $\left(\R^X, G\right)$ to $\left(\R^Y, K\right)$.
\end{definition}
In our context, $F \colon \R^X \to \R^Y$ is non-expansive if for any $\varphi_1, \varphi_2 \in \R^X$:
\[\lVert F(\varphi_1) - F(\varphi_2)\rVert_\infty \le \lVert \varphi_1 - \varphi_2\rVert_\infty.\]
\begin{definition}
A GEO $\left(F, T\right) \colon \left( \R^X, G \right) \to \left( \R^Y, K \right)$ such that $F$ is non-expansive is called a \textbf{Group Equivariant Non-Expansive Operator} (GENEO).
\end{definition}
We say that a GENEO $(F, T)$ is \textit{linear} if the map $F \colon \R^X \to \R^Y$ is a linear map.


% \begin{rem}
% \label{rem:PH}
%     The concept of a GENEO (used in the general setting where finiteness of the
% signal domains is not required; see, e.g., \cite{bergomi2019towards}) provides a
% flexible framework for incorporating different operators into a unified
% setting. A key example in Topological Data Analysis is the computation of
% persistence diagrams, when these are compared using the Hausdorff distance.

% Here we adopt the definition of persistence diagrams in which the diagonal is
% represented as a single point with infinite multiplicity (cf., e.g.,
% \cite{CeEtFr19}). To illustrate this construction, we begin by considering the
% sets 
%     $\Delta=\{(u,u)\in\R\}$,
%     $\Delta^+=\{(u,v)\in\R^2:u < v\}$, 
%     $\Delta^*=\Delta^+\cup\{(u,\infty):u\in\R\}$
%     and $\bar\Delta^*=\Delta^*\cup \{\Delta\}$. We endow $\Delta^*$
%     with the (extended) distance
% \[
% \tilde d(A, B):=\min\left\{\max\left\{|x_A-x_B|,|y_A-y_B|\right\},\max\left\{\frac{|x_A-y_A|}{2},\frac{|x_B-y_B|}{2}\right\}\right\},
% \]
% where $\infty-\infty=0$, 
% $\infty-v=\infty$ for any $v\in\R$, and
% $\frac{\infty}{2}=\infty$. Then, we endow 
%     $\bar\Delta^*$ with the (extended) distance \todo{F: cambierei notazione per essere più simili a PH. $A = (x_A, y_A)$ lo farei diventare $x_1 = (b_1, d_1)$ e analogo con $x_2$. P/N: Ok da parte mia.}
%     \[
% d(A, B) :=
% \begin{cases}
% \tilde d(A, B) & \textup{if } A, B\in \Delta^*,\\
% \frac{y_A-x_A}{2} & \textup{if } A\in \Delta^*, B=\Delta,\\
% \frac{y_B-x_B}{2} & \textup{if } A= \Delta, B\in\Delta^*,\\
% 0 & \textup{if } A=B=\Delta.\\
% \end{cases}
% \]
%     We observe that each nonempty compact subset $C$ of $\bar\Delta^*$ can be represented by
% the function
% $d_C \colon \bar\Delta^* \to \mathbb{R} \cup \{\infty\}$ defined by
% $d_C(p) = \min_{q \in C} d(p,q)$. Analogously to the metric case, one can define the Hausdorff distance $d_H$ associated with the pseudo-metric space $(\bar\Delta^*,d)$.
% Given two nonempty compact subsets $C_1,C_2$ of $\bar\Delta^*$, one can verify that $d_H(C_1,C_2) = \lVert d_{C_1} - d_{C_2} \rVert_\infty$.

% Let now $X$ be a compact metric space and let $\Phi$ the set of all continuous functions from $X$ to $\R$.
% We denote by $\mathrm{Homeo}(X)$ the group of all homeomorphisms of $X$.
% Given $\varphi \in \Phi$ and a fixed homological degree $k$, we denote by
% $\overline{\mathrm{dgm}_k(\varphi)}$ the support of the associated persistence
% diagram in degree $k$, that is, the set of all points of $\bar\Delta^*$ that appear with positive multiplicity
% in the persistence diagram $\mathrm{dgm}_k(\varphi)$ of $\varphi$ in degree $k$. We recall that the support of each persistence diagram is a compact subset of $\bar\Delta^*$ (see~\cite{frosini2025matching}). Denoting the set of persistence diagrams in degree $k$ as $\mathrm{DGM}_k := \{ \mathrm{dgm}_k(\varphi) \mid \varphi \in \Phi \}$, we equip $\mathrm{DGM}_k$ with the pseudo-metric $d_H(\overline{\mathrm{dgm}_k(\varphi_1)},\overline{\mathrm{dgm}_k(\varphi_2)})$. It is easy to check that 
% \[
% (\mathrm{PD}_k,\mathbf{0}) \colon (\Phi, \mathrm{Homeo}(X)) \to (\mathrm{DGM}_k, \{\id\})
% \]
% is a GENEO, where
% $\mathrm{PD}_k$ is the map $\varphi \mapsto \mathrm{dgm}_k(\varphi)$, and $\mathbf{0}$ is
% the trivial homomorphism that maps each homeomorphism of $X$ to the identity
% $\id \colon \mathrm{DGM}_k \to \mathrm{DGM}_k$.
% It is interesting to observe that, in this case, the equivariance of $(\mathrm{PD}_k,\mathbf{0})$ is
% equivalent to the invariance of persistence diagrams with respect to
% reparameterizations of the space $X$, while the non-expansivity of $(\mathrm{PD}_k,\mathbf{0})$ is
% equivalent to the stability of persistence diagrams under perturbations of the
% function $\varphi$ with respect to the max norm.
% \end{rem}

\begin{rem}
\label{rem:PH}
The concept of a GENEO (used in the general setting where finiteness of the signal domains is not required; see, e.g., \cite{bergomi2019towards}) provides a flexible framework for incorporating different operators into a unified setting. A key example in Topological Data Analysis is the computation of persistence diagrams. Let $X$ be a compact metric space, $\Phi$ the set of all continuous functions from $X$ to $\R$, and $\mathrm{Homeo}(X)$ the group of all homeomorphisms of $X$. For a fixed homological degree $k$, 
we consider
%we denote by $\overline{\mathrm{dgm}_k(\varphi)}$ 
the support of the persistence diagram of $\varphi$ in degree $k$, which is a compact subset of the extended half-plane $\bar\Delta^*$, with respect to a suitable metric (see~\cite{frosini2025matching} for details). 
Assuming that the persistence diagram under consideration has no multiple points, we observe that it can be represented both by its support and by the distance function from its support, without any loss of information. In this way, the set
$\mathrm{DGM}_k := \{ \mathrm{dgm}_k(\varphi) \mid \varphi \in \Phi \}$
can be interpreted both as a collection of compact sets and as a space of functions.
Equipping it 
% $\mathrm{DGM}_k := \{ \mathrm{dgm}_k(\varphi) \mid \varphi \in \Phi \}$ 
with the Hausdorff distance $d_H$, one can verify that 
$(\mathrm{PD}_k,\mathbf{0}) \colon (\Phi, \mathrm{Homeo}(X)) \to (\mathrm{DGM}_k, \{\id\})$
is a GENEO, where $\mathrm{PD}_k \colon \varphi \mapsto \mathrm{dgm}_k(\varphi)$ and $\mathbf{0}$ is the trivial homomorphism. The equivariance of $(\mathrm{PD}_k,\mathbf{0})$ corresponds to the invariance of persistence diagrams under reparameterizations of $X$, while non-expansivity corresponds to stability under perturbations of $\varphi$ in the max norm.
\end{rem}

\subsection{Decomposition of stochastic matrices}
This section is devoted to recalling some results on stochastic matrices, mainly adapted from \cite{doi:10.1080/03081087.2020.1733461}. Let $\mathcal{M}_{m \times n}$ be the set of all $m \times n$ real-valued matrices. Moreover, with the symbol $\underline{\ell}$ we denote the set $\{1, \dots, \ell\}$ for every $\ell \in \N \setminus \{0\}$.

\begin{definition}
A matrix $A = (a_{ij}) \in \mathcal{M}_{m \times n}$ is \textbf{(right) stochastic} if $a_{ij} \ge 0$ for all $(i,j) \in \underline{m} \times \underline{n}$ and $\sum_{j=1}^n a_{ij}=1$ for all $i \in \underline{m}$.
\end{definition}

A $\{0,1\}$-matrix is a matrix $A=(a_{ij}) \in \mathcal{M}_{m \times n}$, such that $a_{ij} \in \{0,1\}$ for all $(i,j) \in \underline{m} \times \underline{n}$. We shall refer to the $\{0,1\}$-matrices in $\mathcal{M}_{m \times n}$ with exactly one $1$ in each row as \textbf{rectangular (row) permutation matrices} and $\mathcal{RP}_{m \times n}$ is the set of all rectangular permutation matrices of dimension $m \times n$. 
From now on, we will drop the right and row dependence from the definition of stochastic and rectangular permutation matrices, since it is understood. It is a well-known fact that the set of stochastic matrices is a convex set. Before proceeding we recall Theorem $1$ of \cite{doi:10.1080/03081087.2020.1733461}:

%We recall that an extreme point for a convex set is a point that does not lie in any open line segment joining two points in the set. With this definition, Theorem $1$ of \cite{doi:10.1080/03081087.2020.1733461} can be restated as the following.

\begin{theorem}
\label{thmrectangular}
Every $m \times n$ stochastic matrix can be expressed as a convex combination of $m \times n$ rectangular permutation matrices.
\end{theorem}

\begin{rem}\label{notunique}
    In general, the convex combination stated in Theorem~\ref{thmrectangular} is not unique. As an example, let us consider the following stochastic matrix:
    \[ B = 
    \begin{pmatrix}
        1/2 & 0 & 1/2 \\
        1/3 & 1/3 & 1/3
    \end{pmatrix}.
    \]
    %{\color{red}\sout{To keep the notation simple, we refer to the rectangular permutation matrix with $1$ as $i$-th element of the first row and $j$-th element of the second row as $R_{i,j}$.}}
    To keep the notation simple, we refer to the $2\times 3$ rectangular permutation matrix with entries $1$ in positions $(1,j_1)$ and $(2,j_2)$ as $R_{j_1,j_2}$. Using this notation, $B$ can be expressed as the following convex combinations:
    \[
        B = \frac{1}{12}R_{1,1} + \frac{5}{24}R_{1,2} + \frac{5}{24}R_{1,3} + \frac{1}{4}R_{3, 1} + \frac{1}{8}R_{3, 2} + \frac{1}{8}R_{3,3}
    \]
    and
    \[
        B = \frac{5}{24}R_{1,1} + \frac{1}{12}R_{1,2} + \frac{5}{24}R_{1,3} + \frac{1}{8}R_{3, 1} + \frac{1}{4}R_{3, 2} + \frac{1}{8}R_{3,3}.
    \]
\end{rem}

\section{Building linear GEOs via $T$-permutant measures}
\label{Building}

\sectionmark{Building linear GEOs}
We consider the finite sets $X = \left\{ x_1, \dots, x_n \right\} = \left\{ x_j \right\}_{j=1}^n$ and $Y = \left\{ y_1, \dots, y_m \right\} = \left\{ y_i \right\}_{i=1}^m$ and the relative function spaces $\R^X$ and $\R^Y$.
We fix a group homomorphism $T \colon G \to K$, where $G \subseteq \Aut(X)$ and $K \subseteq \Aut(Y)$. We denote with $X^Y$ the set of all functions from $Y$ to $X$. For each $h \in X^Y$, we consider the following action of $G$ on $X^Y$:
\begin{equation*}
\label{eq:action}
    \alpha_T \colon G \times X^Y \to X^Y, \quad \left( g, h \right) \mapsto g h T(g^{-1}).
\end{equation*}

In particular, we are interested in the orbits $G(h)$ of $h$ under the action $\alpha_T$, for $h \in X^Y$.

\begin{definition}
    A finite signed measure $\mu$ on the power set of $X^Y$ is called a \textbf{(generalized) $T$-permutant measure} if
    %each subset $H$ of $X^Y$ is measurable and
    $\mu$ is invariant under the action $\alpha_T$ of $G$, i.e. $\mu(H) = \mu(gHT(g^{-1}))$ for every $g \in G$. Equivalently, we can say that a finite signed measure $\mu$ on the power set of $X^Y$ is a generalized $T$-permutant measure if
    $\mu(\left\{h\right\}) = \mu(\left\{ghT(g^{-1})\right\})$ for every $g \in G$.
\end{definition}

% \begin{es}[\cite{FaFeFr2023}]\label{oldexample}
% \textcolor{blue}{Francesco: Io metterei $\ell$, visto che come a inizio capitolo è un intero positivo. Nicola metterebbe $d$, visto che è scorrelato dall'uso che ne facciamo sopra. Frosini si erga a giudice.}
% \patrizio{``E dunque, non restando altro lume de la ragione né consiglio d’omini savi, ricorreremo a lo tremendo e infallibile judizio di Dio onnipotente. Sian tratti innanzi li due prodi campioni, armati di cuore, braccio e destino, ché in campo aperto e sotto il guardo de li santi intercessori si portino in singolar tenzone, ciascuno ne lo nome e vessillo de la propria proposta.
% E sappia il mondo tutto che lo detto de lo campione vincitore, per grazia divina e per mezzo de lo brando, sarà accolto, sancito e tenuto per legge, senza pianto né recriminazione. Così volle lo Cielo e così sia.''}
%     Let $X, Y$ be two non-empty finite sets, with $Y \subseteq X$. We fix the equivariance group $G$ on $X$ as the group of all permutations of $X$ that preserve $Y$, and the equivariance group $K$ on $Y$ is the group of all permutations of $Y$. Let us consider the homomorphism $T \colon G \to K$ that takes each permutation of $X$ to its restriction to $Y$. For any positive integer $\ell$, we define $H_\ell = \left\{ h \colon Y \to X \text{ such that } \left| \im(h) \right| = \ell \right\}$, where $\left| \cdot \right|$ denotes the cardinality. We define the finite signed measure $\mu$ as such: given $h \in H_\ell, \mu(h) := \frac{1}{\ell\left| H_\ell \right|}$. Then $\mu$ is a generalized $T$-permutant measure.
% \end{es}

% \patrizio{Devo dire che, a posteriori, l'esempio \ref{oldexample} mi sembra poco chiaro e istruttivo. Lo cambierei nei due che seguono.}

\begin{es}\label{newexample1}
Set $X=\underline{m}\times\underline{n}$ and $Y=\underline{m}$.
We fix $G$ to be the group of all permutations $g$ of $X$ that can be written as
$g=(k,k')$,
where $k$ is a permutation of $Y$ and $k'$ is a permutation of $Z=\underline{n}$, so that 
$g(i,j)=(k(i),k'(j))$.
Then we choose $K$ to be the group of all permutations of $Y$.
The set $X$ can be identified with the set of all cells of an $m\times n$ matrix, while $Y$ represents the set of its rows. Therefore, each $h\in X^{Y}$ can be seen as a map that associates to each row of an $m\times n$ matrix $M$ a cell of $M$.
Consider the homomorphism $T\colon G\to K$
defined by $T(k,k')=k$.
Define a finite signed measure $\mu$ on $X^{Y}$ as follows: for $h\in X^{Y}$, set $\mu(h)=|\operatorname{Im}(h)|$. Then $\mu$ is a generalized $T$-permutant measure.
\end{es}

\begin{es}\label{newexample2}
Using the same notation as in Example~\ref{newexample1}, define a finite signed measure
$\mu$ on $X^{Y}$ as follows: for $h \in X^{Y}$, set
$\mu(h) = 1$ if $h$ maps each row of $M$ into a cell of that row,
and $\mu(h) = 0$ otherwise.
%\todo{F: direi il contrario, se $h$ mappa una cella nella sua rispettiva riga ($h \colon Y \to X$). P: Non capisco: Gli elementi di $Y$ sono le righe, quelli di $X$ le celle.}. 
Then $\mu$ is a generalized $T$-permutant
measure.
% Note that for each $h \in X^{Y}$, the set
% $\{\, g h T(g^{-1}) : g \in G \,\}$ consists of all maps $\tilde h : Y \to X$
% such that, for every index $i$, both $\tilde h(i)$ and $h(i)$ belong to the same
% row of $M$.
\end{es}


Moreover, with a slight abuse of notation, we denote with $\mu(h)$ the signed measure of the singleton $\left\{ h \right\}$ for each $h \in X^Y$. Thanks to the concept of generalized $T$-permutant measure, we are able to define a linear GEO associated with it. In fact, we can state the following.

\begin{proposition}
    \label{prop:GEOperm}
    Let $\mu$ be a generalized $T$-permutant measure. We define the operator $F_\mu \colon \R^X \to \R^Y$ associated with the measure $\mu$ by setting
    \[
    F_\mu \left( \p \right) := \sum_{h \in X^Y} \p h \ \mu(h)
    \]
    for every $\p \in \R^X$. Then $(F_\mu,T) \colon (\R^X,G) \to (\R^Y,K)$ is a linear GEO.
\end{proposition}
\begin{proof}
    First, we show that $F_\mu$ is $T$-equivariant. It holds that:
    \begin{align*}
        F_\mu \left( \p g \right) & = \sum_{h \in X^Y} \p g h \ \mu(h) \\
        & = \sum_{h \in X^Y} \p g h T (g^{-1}) T (g) \ \mu\left(g h T(g^{-1})\right) \\
        & = \sum_{f \in X^Y} \p f T(g)\ \mu(f) \\
        & = F_\mu \left( \p \right) T(g),
    \end{align*}
    since $\mu(h) = \mu \left( ghT(g^{-1}) \right)$ and the map $h \mapsto f := g h T(g^{-1})$ is a bijection from $X^Y$ to $X^Y$. The linearity of $F_\mu$ directly comes from the definition.
    % \textcolor{red}{N: By definition $F$ is linear. composizione per h è lineare.} Now we show that $F$ is linear. Let $\p, \psi \in \R^X$ and let $\lambda \in \R$. It holds that:
    % \begin{align*}
    %     F \left( \p + \lambda \psi \right) & = \sum_{h \in X^Y} \left( \p + \lambda \psi \right) h \mu(h) \\
    %     & = \sum_{h \in X^Y} \p h \mu(h) + \sum_{h \in X^Y} \lambda\psi h \mu(h) \\
    %     & = F(\p) + \lambda F(\psi).
    % \end{align*}
    % Therefore, $F$ is a linear GEO.
\end{proof}

\section{A representation theorem for linear GEOs via $T$-permutant measures} 
\label{Representation}

Proposition~\ref{prop:GEOperm} provides a first link between generalized $T$-permutant measures and GEOs, and its definition and proof was straightforward. In contrast, proving the other verse of the representation theorem, i.e. associating a generalized $T$-permutant measure with a GEO, is more difficult. In any case, we are able as of now to state the result that we want to prove.

% \begin{theorem}
%     \label{maint}
%     \textcolor{purple}{N: Correggere "surjective" con le nuove ipotesi. E correggere il testo che risulta desueto.} 
%     Assume that $G \subseteq \mathrm{Aut}(X), K \subseteq \mathrm{Aut}(Y)$, $K$ transitively acts on the finite set $Y$, $T \colon G \to K$ is a surjective homomorphism and $F$ is a map from $\mathbb{R}^X$ to $\mathbb{R}^Y$. The map $F$ is a linear group equivariant operator from $(\mathbb{R}^X, G)$ to $(\mathbb{R}^Y, K)$ with respect to the homomorphism $T$ if and only if a generalized permutant measure $\mu$ exists such that $F(\varphi) = \sum_{h \in X^Y} \varphi h \mu(h)$ for every $\varphi \in \mathbb{R}^X$.
% \end{theorem}

\begin{theorem}
    \label{maint} 
    Assume that $G \subseteq \mathrm{Aut}(X), K \subseteq \mathrm{Aut}(Y)$, and $T \colon G \to K$ is a group homomorphism. Moreover, suppose that $T(G)$ transitively acts on the finite set $Y$, and $F$ is a map from $\mathbb{R}^X$ to $\mathbb{R}^Y$. The map $(F,T)$ is a linear GEO from $(\mathbb{R}^X, G)$ to $(\mathbb{R}^Y, K)$ if and only if a $T$-permutant measure $\mu$ exists such that $F(\varphi) = \sum_{h \in X^Y} \varphi h\  \mu(h)$ for every $\varphi \in \mathbb{R}^X$.
\end{theorem}
%\textcolor{red}{N: Suriettività e transività sono entrambe da richiedere? Transitivity does not imply that $T$ is surjective. Indeed, the trivial homomorphism can be always considered.}
%\begin{rem}
%In the non-homogeneous setting, the action of $T(G)$ on $Y$ decomposes $Y$ into several distinct orbits.  
%This situation may look more complicated at first, yet it does not introduce any real difficulty: one can simply project $F$ onto each orbit and study its behavior separately.  
%Each component then falls within the framework of Theorem~\ref{teo:linear}, so the general case can be regarded as a combination of the homogeneous ones.
%\end{rem}

\begin{rem}
% \sout{In the non-homogeneous setting,} 
% \pat{Di solito per spazio omogeneo si intende una varietà e questo non è il nostro caso}
In the case where the action of $T(G)$ on $Y$ is not transitive, 
this action decomposes $Y$ into several distinct orbits $Y_r$.
Even in this case, however, we can take advantage of Theorem~\ref{maint}, applying it to study each restriction $F(\p)_{\vert Y_r}:Y_r\to\R$ as the orbit $Y_r$ varies. This is possible since the action of $T(G)$ on $Y_r$ is clearly transitive.
\end{rem}

Throughout this paper, we use $\varphi$ to denote a function
$\varphi : X \to \mathbb{R}$, and $[\varphi]$ to denote its associated column
vector representation with respect to the standard basis of $\mathbb{R}^X$.
Similarly, we use $F$ to denote a linear map
$F : \mathbb{R}^X \to \mathbb{R}^Y$, and $[F]$ to denote its associated matrix
with respect to the standard bases $\mathcal{B}_X=([{\mathds{1}}_{x_1}],\ldots,[{\mathds{1}}_{x_n}])$ of $\R^X$ and $\mathcal{B}_Y=([{\mathds{1}}_{y_1}],\ldots [{\mathds{1}}_{y_m}])$ of $\R^Y$. The remainder of this section is devoted to the proof of Theorem~\ref{maint}.
We split the proof into several substeps, both to improve readability and to
present each step as a separate result. We stress the fact that one verse of the proof has already been proved by Proposition~\ref{prop:GEOperm}. Let us assume that $(F,T) \colon (\R^X,G) \to (\R^Y,K)$ is a linear GEO. Moreover, let $B=(b_{ij})$ be the matrix
% \todo{F: le definizioni di $[F]$ e $B$ sono praticamente attaccate, quindi io scriverei due parole sul fatto che coincidono ma usiamo $[F]$ quando vogliamo dire qualcosa sull'operatore, mentre usiamo $B$ quando la trattiamo puramente come matrice.}} 
associated with $F$ with respect to the standard bases of $\R^X$ and $\R^Y$. The meanings of the symbols \(B\) and \([F]\) are therefore the same, but we will use the latter when we wish to emphasize that it is the matrix associated with the chosen operator \(F\). Given a natural number \(\ell \in \mathbb{N} \setminus \{0\}\) and a set \(Z\) with \(\lvert Z \rvert = \ell\), for every permutation \(p \colon Z \to Z\) we denote by \(\sigma_p \colon \underline{\ell} \to \underline{\ell}\) the function defined by setting \(\sigma_p(j) = i\) if and only if \(p(z_j) = z_i\). We observe that \(\sigma_{p^{-1}} = \sigma_p^{-1}\) and that
$\mathds{1}_{x_j} g = \mathds{1}_{g^{-1}(x_j)}$.
A similar identity holds for \(\mathds{1}_{y_i} T(g)\).

\begin{lemma}\label{rijhm}
For any $g \in G$, we have that $b_{ij}=b_{\sigma_{T(g)}(i)\sigma_g(j)}$ for every $(i,j) \in \underline{m} \times \underline{n}$.
\end{lemma}

\begin{proof}
Let us choose a function ${\mathds{1}}_{x_j}$ and a permutation $g \in G$.
By equivariance we have that
\[{F}({\mathds{1}}_{x_j}g )= {F}({\mathds{1}}_{x_j})T(g).\]
The left-hand side of the equation can be rewritten as:
\[
{F}({\mathds{1}}_{x_j}g )= {F}({\mathds{1}}_{g^{-1}(x_j)} )= \sum_{i=1}^{m}b_{i\sigma_g^{-1}(j)}{\mathds{1}}_{y_i}.
\]
On the right-hand side, we get
\begin{align*}
{F}({\mathds{1}}_{x_j})T(g)&=\left(\sum_{i=1}^{m}b_{ij}{\mathds{1}}_{y_i}\right) T(g)
\\&= \sum_{i=1}^{m}b_{ij}({\mathds{1}}_{y_i}T(g))
\\&= \sum_{i=1}^{m}b_{ij}({\mathds{1}}_{T(g^{-1})(y_i)})
\\&=\sum_{s=1}^{m}b_{\sigma_{T(g)}(s)j}{\mathds{1}}_{y_s},
\end{align*}
by setting $y_s=T(g^{-1})(y_i)$,
{\color{black} and hence $i=\sigma_{T(g)}(s)$.}
Therefore, we obtain the following equation: 
\[
\sum_{i=1}^{m}b_{i\sigma_g^{-1}(j)}{\mathds{1}}_{y_i}=\sum_{s=1}^{m}b_{\sigma_{T(g)}(s)j}{\mathds{1}}_{y_s}.
\]
This immediately implies that $b_{i\sigma_g^{-1}(j)}=b_{\sigma_{T(g)}(i)j}$, for any $i \in \underline{m}$.
Since this equality holds for
any $j \in \underline{n}$ and any $g \in G$, we have that $b_{ij}=b_{\sigma_{T(g)}(i)\sigma_g(j)}$ for every $(i,j) \in \underline{m} \times \underline{n}$ and every $g\in G$.
\end{proof}

For the rest of the section, we will assume that $T(G)$ transitively acts on $Y$. 
% and $T$ is surjective. 
% For the rest of the section, we will assume that $K$ transitively acts on $Y$ and $T$ is surjective. \textcolor{purple}{Ci serve veramente la suriettività di $T$?}
%We note that \textcolor{purple}{$T$-equivariance} and the surjectivity of $T$ implies that $m \ge n$ \textcolor{purple}{perchè è implicata 'sta cosa? Quando la usiamo?} \textcolor{blue}{F: la usiamo implicitamente dopo la dimostrazione del lemma 2.10, se non fosse $m \ge n$, h non sarebbe una funzione. Per la dimostrazione è vero, dobbiamo provarlo}. 
These assumptions allow us to prove the following result.

\begin{lemma}\label{lempermutationrow}
If the action of $T(G) \subseteq K$ on $Y$ is transitive, an $n$-tuple of real numbers $\beta=(\beta_1,\ldots,\beta_n)$ exists such that each row of ${B}$ can be obtained by permuting $\beta$.
\end{lemma}
% \begin{proof}
% Since $K$ acts transitively, for every $i \in [n]$ there exists
% $h_{i \mapsto 1}\in K$ such that $h_{i\mapsto1}(y_i)=y_1$.
% Considering the $\overline \imath$-th row of $B$ and $h_{\overline \imath \mapsto 1} \in K$, there exists $\overline g \in G$ such that $T(\overline g)=h_{\overline \imath \mapsto 1}$ because $T$ is surjective. Hence, by Lemma \ref{rijhm}, we have that $b_{\overline \imath j}
% =b_{\sigma_{T (\overline g)}(\overline \imath) \sigma_{\overline g}(j)}
% =b_{\sigma_{h_{\overline\imath \mapsto 1}}(\overline \imath) \sigma_{\overline g}(j)}
% =b_{1\sigma_{\overline g}(j)}$, for any $j \in [m]$.
% Since $\sigma_{\overline g}$ is a permutation, the $\overline\imath$-th row is a permutation of the first row.
% \end{proof}
\begin{proof}
Since $T(G)$ acts transitively, for every $i \in \underline{m}$ there exists
$k_{i \mapsto 1}\in T(G)$ such that $k_{i\mapsto1}(y_i)=y_1$.
Considering the $\overline \imath$-th row of $B$ and $k_{\overline \imath \mapsto 1} \in T(G)$, there exists $\overline g \in G$ such that $T(\overline g)=k_{\overline \imath \mapsto 1}$. Hence, by Lemma~\ref{rijhm}, we have that $b_{\overline \imath j}
=b_{\sigma_{T (\overline g)}(\overline \imath) \sigma_{\overline g}(j)}
=b_{\sigma_{k_{\overline\imath \mapsto 1}}(\overline \imath) \sigma_{\overline g}(j)}
=b_{1\sigma_{\overline g}(j)}$, for any $j \in \underline{n}$.
Since $\sigma_{\overline g}$ is a permutation, the $\overline\imath$-th row is a permutation of the first row.
\end{proof}

% \begin{proof}
% Since $T(G)$ acts transitively, for every $i \in \underline{m}$ there exists
% $h_{i \mapsto 1}\in T(G)$ such that $h_{i\mapsto1}(y_i)=y_1$.
% Considering the $\overline \imath$-th row of $B$ and $h_{\overline \imath \mapsto 1} \in T(G)$, there exists $\overline g \in G$ such that $T(\overline g)=h_{\overline \imath \mapsto 1}$. Hence, by Lemma~\ref{rijhm}, we have that $b_{\overline \imath j}
% =b_{\sigma_{T (\overline g)}(\overline \imath) \sigma_{\overline g}(j)}
% =b_{\sigma_{h_{\overline\imath \mapsto 1}}(\overline \imath) \sigma_{\overline g}(j)}
% =b_{1\sigma_{\overline g}(j)}$, for any $j \in \underline{n}$.
% Since $\sigma_{\overline g}$ is a permutation, the $\overline\imath$-th row is a permutation of the first row.
% \end{proof}

In the following, it is helpful to split $F$ in its positive and negative parts. That is, let us consider the linear maps ${F^\oplus},{F^\ominus}:\R^X\to\R^Y$ defined by setting ${F^\oplus}(\mathds{1}_{x_j}) := \sum_{i=1}^{m}\max \left\{b_{ij},0\right\} \mathds{1}_{y_i}$ and ${F^\ominus}(\mathds{1}_{x_j}) := \sum_{i=1}^{m}\max \left\{-b_{ij},0\right\} \mathds{1}_{y_i}$ for every index $j\in \underline{n}$. 
{\color{black} By applying Lemma~\ref{rijhm},} one can easily check that the following properties hold:

\begin{enumerate}
  \item ${F^\oplus},{F^\ominus}$ are $T$-equivariant linear maps;
  \item The matrices $B^\oplus$ and $B^\ominus$ associated with ${F^\oplus}$ and ${F^\ominus}$ with respect to the bases $\mathcal{B}_X$ and $\mathcal{B}_Y$ are $B^\oplus =\left(b^\oplus_{ij}\right) = \left( \max\left\{b_{ij},0\right\}\right)$ and $B^\ominus = \left(b^\ominus_{ij}\right)=\left(\max\left\{-b_{ij},0\right\}\right)$, respectively. In particular, $B^\oplus$ and $B^\ominus$ are non-negative matrices;
  \item $b_{ij}^\oplus \neq 0 \implies b_{ij}^\ominus = 0$, and $b_{ij}^\ominus \neq 0 \implies b_{ij}^\oplus = 0$ for any pair of indices $i,j$;
  \item $F=F^\oplus - F^\ominus$ and $B={B^\oplus}-{B^\ominus}$;
    \item \label{property5} Lemma~\ref{lempermutationrow} and the definitions of $B^\oplus, B^\ominus$ imply that two $n$-tuples of positive real numbers $\beta^\oplus=\left(\beta^\oplus_1,\dots,\beta^\oplus_{n}\right)$, $\beta^\ominus=\left(\beta^\ominus_1,\dots,\beta^\ominus_{n}\right)$ exist such that each row of ${B^\oplus}$ can be obtained by permuting $\beta^\oplus$, and each row of ${B^\ominus}$ can be obtained by permuting $\beta^\ominus$.
\end{enumerate}

The next step in order to prove Theorem~\ref{maint} is to express $F^\oplus$ and $F^\ominus$ as weighted sums of $\varphi h$, for $h \in X^Y$. First, we need to establish a connection between the elements of $X^Y$ and the rectangular permutation matrices. Each function $h \colon Y \to X$ can be associated with a $m \times n$ rectangular permutation matrix $R(h)=(r_{ij}(h))$ defined by setting $r_{ij}(h)=1$ if $h(y_i)=x_j$, and $r_{ij}(h)=0$ otherwise. In the case $X = Y$ and $h \colon X \to X$ is a permutation, we denote as $P(h)$ such square matrix, which is the usual permutation matrix. 
%Furthermore, one can prove that
Hence, there is a bijection between $X^Y$ and the set $\mathcal{RP}_{m \times n}$
of all $m \times n$ rectangular permutation matrices, given by $h \mapsto R(h)$.


\begin{rem}\label{grhughaepj}
Assume that $h$ is a function from $Y$ to $X$ and $\mathcal{R}_h$ denotes the linear operator that sends each function $\varphi$ in $\R^X$ to $\varphi h$ in $\R^Y$. One could easily check that $R(h)$ is the matrix associated with the operator $\mathcal{R}_h$ with respect to the bases $\mathcal{B}_X$ for $\R^X$ and $\mathcal{B}_Y$ for $\R^Y$. 
Moreover, considering two functions ${f_1} \colon X \to Y$, ${f_2} \colon Y \to Z$, and a function $\p\in \R^Z$, we know that $\mathcal{R}_{{f_2} {f_1}}(\varphi)= \varphi ({f_2} {f_1}) = (\varphi {f_2}) {f_1} = \mathcal{R}_{f_1} \mathcal{R}_{f_2}(\varphi)$. Hence, $ R({f_2}{f_1})=R({f_1}) R({f_2})$.
\end{rem}

\begin{proposition}
\label{prop:sum}
    For every $h \in X^Y$ there exist two non-negative real numbers $c^\oplus(h), c^\ominus(h)$ such that
    \begin{align*}
    F^\oplus(\varphi) & = \sum_{h \in X^Y}c^\oplus(h) \varphi h, \\
    F^\ominus(\varphi) & = \sum_{h \in X^Y}c^\ominus(h) \varphi h,
    \end{align*}
    for every $\varphi \in \R^X$.
    Therefore  
\begin{align*}
    B^\oplus & = \sum_{h \in X^Y} c^\oplus(h)\, R(h), \\
    B^\ominus & = \sum_{h \in X^Y} c^\ominus(h)\, R(h).
\end{align*}
\end{proposition}

\begin{proof}
    Let us start by considering the statement concerning $c^\oplus$ and $F^\oplus$. \textcolor{black}{If $F^\oplus$ is the zero operator, the statement follows immediately.} {\color{black} Otherwise, applying Lemma~\ref{lempermutationrow} and setting $\lVert \beta^\oplus\rVert_1:= \sum_{i=1}^n \beta^\oplus_i\neq 0$}, we note that the matrix $\frac{1}{\lVert \beta^\oplus\rVert_1} B^\oplus$ is stochastic. 
    Hence, Remark~\ref{grhughaepj} and Theorem~\ref{thmrectangular} imply that
    % \patrizio{[Non mi piace confondere la funzione $F^\oplus(\varphi)$ con un vettore, nonostante che a pagina 6 abbiamo messo le mani avanti dicendo che lo avremmo fatto. Preferirei scrivere 
    % $[F^\oplus(\varphi)]$ per indicare il vettore delle coordinate della funzione $F^\oplus(\varphi)$.]}
    \begin{align*}
    [F^\oplus(\varphi)] & = \lVert \beta^\oplus\rVert_1 \left( \frac{1}{\lVert \beta^\oplus\rVert_1}B^\oplus [\varphi]\right) \\
    & = \lVert \beta^\oplus\rVert_1 \sum_{h \in X^Y}\gamma^\oplus (h)R(h)[\varphi] \\
    & = \sum_{h \in X^Y} c^\oplus(h) [\mathcal{R}_h (\varphi)] \\
    & = \left[\sum_{h \in X^Y} c^\oplus(h) \varphi h\right],
    \end{align*}
    where $\gamma^\oplus(h)$ is a nonnegative real number, $\sum_{h \in X^Y}\gamma^\oplus(h) = 1$, and $c^\oplus(h)= \lVert\beta^\oplus\rVert_1 \gamma^\oplus(h)\ge 0$ for any $h \in X^Y$.
    {\color{black} Therefore, $F^\oplus(\varphi)=\sum_{h \in X^Y} c^\oplus(h) \varphi h$.}
    \textcolor{black}{
    Moreover, it follows that
    \begin{align*}
    B^\oplus & = \lVert \beta^\oplus\rVert_1 \left( \frac{1}{\lVert \beta^\oplus\rVert_1}B^\oplus \right) \\
    & = \lVert \beta^\oplus\rVert_1 \sum_{h \in X^Y}\gamma^\oplus (h)R(h) \\
    & = \sum_{h \in X^Y} c^\oplus(h) R(h).
    \end{align*}
    }
    The proof of the statements concerning $c^\ominus$, $F^\ominus$, and $B^\ominus$ is analogous.
\end{proof}

\begin{rem}\label{gammanotunique}
It is important to observe that, due to Remark~\ref{notunique}, the choice of the
coefficients $\gamma^\oplus(h)$ and $\gamma^\ominus(h)$ in the proof of
Proposition~\ref{prop:sum} is not unique in general.
\end{rem}

We recall that our aim is to associate a generalized $T$-permutant measure to a
GEO $(F,T)$, and we split this task into its positive and negative parts
$F^\oplus$ and $F^\ominus$. At this stage, the function $h \mapsto c^\oplus(h)$
is not a generalized $T$-permutant measure. In order to obtain a generalized
$T$-permutant measure, we need to average it over the orbits of $h$ under the
action of $G$. First, we show that this averaging procedure is well defined.


\begin{proposition}\label{propc}
If $h_1, h_2 \in X^Y$ and there exists $t \in \underline{m}$ such that $h_1(y_t) = h_2(y_t)$, then either $c^\oplus(h_1)=0$, or $c^\ominus(h_2)=0$ or both values are null.
\end{proposition}

\begin{proof}
    Since there exist $t \in \underline{m}$ {\color{black} and $s\in\underline{n}$} such that $x_s=h_1(y_t) = h_2(y_t)$, we can consider $\mathds{1}_{x_s}$, and get that
    % \pat{Anche qui userei $[F^\oplus \left( \mathds{1}_{x_s} \right)] $ al posto di $F^\oplus \left( \mathds{1}_{x_s} \right) $. Inoltre ho sostituito $n$ con $m$ nelle due sommatorie in $i$ che seguono.}
    \begin{align*}
        [F^\oplus \left( \mathds{1}_{x_s} \right)] & = B^\oplus [\mathds{1}_{x_s}] \\
        & = \left[\sum_{i = 1}^m b^\oplus_{is}\mathds{1}_{y_i}\right],
    \end{align*}
    {\color{black} i.e., $F^\oplus \left( \mathds{1}_{x_s} \right)=\sum_{i = 1}^m b^\oplus_{is}\mathds{1}_{y_i}$.}
    Hence, we have that
\begin{align*}
        b^\oplus_{ts} & = \left( \sum_{i = 1}^m b^\oplus_{is}\mathds{1}_{y_i}\right)(y_t)
        \\ & = F^\oplus \left( \mathds{1}_{x_s} \right)(y_t)
        \\& = \sum_{h \in X^Y}c^\oplus(h) \mathds{1}_{x_s}h(y_t)
        \\ & \ge c^\oplus(h_1) \mathds{1}_{x_s}h_1(y_t)
        \\ & = c^\oplus(h_1) \mathds{1}_{x_s}(x_s)
        \\ & = c^\oplus(h_1).
    \end{align*}
    One could similarly check that $b^\ominus_{ts} \ge c^\ominus({h_2})$. Therefore,
    \[
    c^\oplus({h_1})>0 \implies b^\oplus_{ts}>0 \implies b^\ominus_{ts}=0 \implies c^\ominus({h_2})=0.
    \]
    It follows that either $c^\oplus(h_1)=0$, or $c^\ominus(h_2)=0$, or both.
\end{proof}

\begin{corollary}
\label{corc}
For every $h\in X^Y$, either $c^\oplus(h)=0$, or $c^\ominus(h)=0$, or both.
\end{corollary}

\begin{proof}
    Set $h_1=h_2$ in Proposition~\ref{propc}.
\end{proof}

We finally have the prerequisites to define the generalized $T$-permutant measures $\mu^\oplus$ and $\mu^\ominus$ associated with $F^\oplus$ and $F^\ominus$. Each of them is the average of the functions $c^\oplus$ and $c^\ominus$ along the orbit $\mathcal{O}(h)$ of $h$ under the action $\alpha_T$. Formally, we define
\begin{align*}
    \mu^\oplus(h) & := \sum_{f \in \mathcal{O}(h)} \frac{c^\oplus(f)}{\left|\mathcal{O}(f)\right|} = \sum_{f \in \mathcal{O}(h)} \frac{c^\oplus(f)}{\left|\mathcal{O}(h)\right|}, \\
    \mu^\ominus(h) & := \sum_{f \in \mathcal{O}(h)} \frac{c^\ominus(f)}{\left|\mathcal{O}(f)\right|} = \sum_{f \in \mathcal{O}(h)} \frac{c^\ominus(f)}{\left|\mathcal{O}(h)\right|}.
\end{align*}

\begin{proposition}
\label{prop:gpm}
    $\mu^\oplus$ and $\mu^\ominus$ are generalized $T$-permutant measures. As a consequence, the function $\mu = \mu^\oplus - \mu^\ominus$ is also a generalized $T$-permutant measure.
\end{proposition}

\begin{proof}
    The definition of $\mu^\oplus$ immediately implies that $\mu^\oplus(H) = \mu^\oplus\left(gHT(g)^{-1}\right)$ for every $g \in G$ and every subset $H$ of $X^Y$. In other words, $\mu^\oplus$ is a non-negative generalized $T$-permutant measure. Quite analogously, we can prove that $\mu^\ominus$ is a non-negative generalized $T$-permutant measure. Hence, we can conclude that the function $\mu := \mu^\oplus - \mu^\ominus$ is a generalized $T$-permutant measure.
    % \patrizio{
    % Non capisco a che serva qui il Corollario \ref{corc}.
    % Mi sembra che si possa affermare che 
    % $\mu := \mu^\oplus - \mu^\ominus$ è una misura permutante generalizzata a prescindere da quel corollario. Mi pare che il corollario serva solo nel seguito.
    % Per questo motivo, se ricordo bene, in \cite{BoBoBrFrQu23} l'analogo del Corollario \ref{corc} era stato inserito dopo la definizione di $\mu$.}
\end{proof}

Now that we have defined the two generalized $T$-permutant measures
$\mu^\oplus$ and $\mu^\ominus$ associated with the $T$-equivariant linear maps
$F^\oplus$ and $F^\ominus$, we can begin the steps needed to show that suitable
weighted sums of these measures reproduce precisely the two maps. 
Let $G_h:=\{g \in G \mid \alpha_T(g,h)=h\}$ be the stabilizer subgroup of $G$ with respect to $h$, i.e., the subgroup of $G$ containing the elements that fix $h$ by the action. We recall that acting on $h$ by every element of $G$ we obtain each element of the orbit $\mathcal{O}(h)$ exactly $\left|G_h\right|$ times, and the well known relation $\left|G_h\right|\left|\mathcal{O}(h)\right| = \left|G\right|$ (cf. \cite{aschbacher_2000}). We observe that, for $f, h \in X^Y$, if $f \in \mathcal{O}(h)$ then $G_f$ is isomorphic to $G_h$. Finally we are able to prove Proposition~\ref{carperm+-}, which states something very similar to Theorem~\ref{maint} but regarding $F^\oplus$ and $F^\ominus$.

\begin{proposition} 
    \label{carperm+-}
    For any $\varphi \in \R^X$ we have that
    \begin{align*}
        F^\oplus(\varphi) & = \sum_{h \in X^Y} \varphi h \ \mu^\oplus(h), \\
        F^\ominus(\varphi) & = \sum_{h \in X^Y} \varphi h \ \mu^\ominus(h).
    \end{align*}
Therefore,
 \begin{align*}
        B^\oplus & = \sum_{h \in X^Y} \mu^\oplus(h) R(h), \\
        B^\ominus & = \sum_{h \in X^Y} \mu^\ominus(h)R(h).
    \end{align*}
\end{proposition}

\begin{proof}
    Recalling that $\mathcal{R}_g \colon \varphi \mapsto \varphi g$ and $\mathcal{R}_k \colon \psi \mapsto \psi k $ are linear maps for any $g \in G$ and $k \in K$, the $T$-equivariance condition $F^\oplus \mathcal{R}_{g^{-1}} = \mathcal{R}_{T(g^{-1})}F^\oplus$ directly implies that $B^\oplus P(g^{-1}) = P\left(T(g^{-1})\right)B^\oplus$. In particular, we have that $P\left(T(g)\right)B^\oplus P(g)^{-1} = B^\oplus$ for every $g \in G$. Note that $P(g)$ is an $n\times n$ matrix, while $P(T(g))$ is an $m\times m$ matrix. From Proposition~\ref{prop:sum} it follows that
    %{\allowdisplaybreaks
    \begin{align*}
        B^\oplus & = \overbrace{\frac{1}{\left|G\right|}B^\oplus + \dots + \frac{1}{\left|G\right|}B^\oplus}^{\left|G\right| \text{ summands}} \\
        & = \frac{1}{\left|G\right|}\sum_{g \in G}P(T(g))B^\oplus P(g)^{-1} \\
        & = \frac{1}{\left|G\right|}\sum_{g \in G}P(T(g))\left(\sum_{f \in X^Y} c^\oplus(f)R(f)\right)P(g)^{-1} \\
        & = \sum_{f \in X^Y}\sum_{g \in G}\frac{c^\oplus(f)}{\left|G\right|} P(T(g))R(f)P(g)^{-1} \\
        & = \sum_{f \in X^Y}\frac{c^\oplus(f)}{\left|G\right|}\sum_{g \in G} R\left(g^{-1}fT(g)\right) \tag{$\star$} \\
        & = \sum_{f \in X^Y}\frac{c^\oplus(f)}{\left|G\right|}\sum_{g \in G} R\left(gfT(g)^{-1}\right),
    \end{align*}
    %}
    where $(\star)$ follows from Remark~\ref{grhughaepj}, recalling that $P(g)=R(g)$ and $P(T(g))=R(T(g))$. Therefore, we have that
    \begin{align}
        [F^\oplus(\varphi)] & = B^\oplus[\p] \nonumber\\
        & = \sum_{f \in X^Y} \frac{c^\oplus(f)}{\left|G\right|} \sum_{g \in G} R\left(gfT(g)^{-1}\right)[\varphi] \nonumber \\
        & = \left[\sum_{f \in X^Y} \frac{c^\oplus(f)}{\left|G\right|} \sum_{g \in G} \mathcal{R}_{gfT(g)^{-1}}(\varphi)\right] \nonumber\\
        & = \left[\sum_{f \in X^Y} \frac{c^\oplus(f)}{\left|G\right|} \sum_{g \in G} \varphi g f T(g)^{-1}\right],\nonumber
    \end{align}
    i.e., 
    \begin{align}
    \label{eq:conj}
        F^\oplus(\varphi) & = \sum_{f \in X^Y} \frac{c^\oplus(f)}{\left|G\right|} \sum_{g \in G} \varphi g f T(g)^{-1}.
    \end{align}
    Let us now set $\delta(f_1, f_2) = 1$ if $f_1$ and $f_2$ belong to the same orbit under the action of $G$, and $\delta(f_1, f_2) = 0$ otherwise. Moreover, we recall that, by acting on $h$ with every element of $G$, each element of the orbit $\mathcal{O}(h)$ is obtained exactly $|G_h| = \frac{|G|}{|\mathcal{O}(h)|}$ times, and that $f \in \mathcal{O}(h) \iff h \in \mathcal{O}(f)$. Therefore, Equation~(\ref{eq:conj}) implies that
    \begin{align*}
        F^\oplus(\varphi) & = \sum_{f \in X^Y} \frac{c^\oplus(f)}{\left|G\right|} \left|G_f\right| \sum_{h \in \mathcal{O}(f)} \varphi h \\
        & = \sum_{f \in X^Y} \frac{c^\oplus(f)}{\left|G\right|} \left|G_f\right| \sum_{h \in X^Y} \delta(h, f)\varphi h \\
        & = \sum_{h \in X^Y} \left( \sum_{f \in X^Y} \frac{c^\oplus(f)}{\left|G\right|} \left|G_f\right| \delta(h, f) \right) \varphi h \\
        & = \sum_{h \in X^Y} \left( \sum_{f \in X^Y} \frac{c^\oplus(f)}{\left|\mathcal{O}(f)\right|} \delta(h, f) \right) \varphi h \\
        & = \sum_{h \in X^Y} \left( \sum_{f \in \mathcal{O}(h)} \frac{c^\oplus(f)}{\left|\mathcal{O}(f)\right|} \right) \varphi h \\
        & = \sum_{h \in X^Y} \varphi h \ \mu^\oplus(h).
    \end{align*}
    The equality $B^\oplus = \sum_{h \in X^Y} \mu^\oplus(h) R(h)$ follows by observing that 
    $[F^\oplus(\p)]=B^\oplus [\p]$ and
    $\left[\varphi h \ \mu^\oplus(h)\right]=\mu^\oplus(h) R(h)[\p]$ for every $h\in X^Y$ and $\p\in \R^X$. The statements concerning $F^\ominus$ and $B^\ominus$ are proved analogously.
\end{proof}

We can now prove Theorem~\ref{maint}.
\begin{proof}[Proof of Theorem~\ref{maint}]
   Proposition~\ref{prop:GEOperm} already provides one verse of the theorem. From Proposition~\ref{prop:gpm} and Proposition~\ref{carperm+-}, it follows that $\mu$ is a generalized $T$-permutant measure such that $F(\varphi) = \sum_{h \in X^Y} \varphi h \ \mu(h)$ for every $\varphi \in \mathbb{R}^X$.
\end{proof}

% The previous construction makes it possible to introduce the following definition.

% \begin{definition}
% Each measure $\mu$ constructed according to the proof of Theorem~\ref{maint} will be called
% a \textit{measure associated with the GEO $F$}.
% \end{definition}

% % We recall that this section requires that $T(G)$ transitively acts on $Y$. We have just proved a representation theorem for linear GEOs. Now, we are going to extend it to include non-expansivity. Before proceeding, we want to say something more about the relationship between $\mu^\oplus$ and $\mu^\ominus$.

\begin{rem}\label{measurenotunique}
Because of Remark~\ref{gammanotunique}, the measures $\mu^\oplus$, $\mu^\ominus$ and $\mu$ associated with $F$ in the proof of Theorem~\ref{maint} are not unique in general.
\end{rem}

We wish to emphasize here that the statement of our representation theorem for
linear GEOs (Theorem~\ref{maint}) requires that $T(G)$ acts transitively on $Y$. Before extending this result to include non-expansivity, we make an additional remark on the relationship between $\mu^\oplus$ and $\mu^\ominus$.

\begin{proposition}
    \label{prop:singular}
    The measures $\mu^\oplus$ and $\mu^\ominus$ associated with $F$ in the proof of Theorem~\ref{maint} are mutually singular.
\end{proposition}

% \begin{proof}
%     If $\mu^\oplus \equiv 0$ there is nothing to prove. Let us assume that $\mu^\oplus$ is not the null measure and consider $\overline{h} \in X^Y$ such that $\mu^\oplus\left(\overline{h}\right)>0$. By contradiction, we suppose that $\mu^\ominus\left(\overline{h}\right)>0$. Then, by definition of $\mu^\ominus$ there exists $f \in \mathcal{O}\left(\overline{h}\right)$ such that $c^\ominus(f) > 0$. Since $R(f) = (r(f)_{ij})$ is a rectangular permutation matrix, there is an index $(\overline{\imath}, \overline{\jmath})$ such that $r(f)_{\overline{\imath}, \overline{\jmath}} = 1$. Hence, by Proposition~\ref{prop:sum}, it holds that $b^\ominus_{\overline{\imath}, \overline{\jmath}} = \sum_{h \in X^Y} r(h)_{\overline{\imath}, \overline{\jmath}} c^\ominus(h) > 0$. \textcolor{blue}{F: A me pare che l'unica volta che usiamo la proprietà 3 sopra (che non ha nome) sia qui, per cui l'enuncerei direttamente qui e basta, oppure diamo un nome a quelle 5 proprietà (es. Properties) e qui la richiamiamo. Cambiare la frase dopo a seconda di quello che vogliamo fare.}Since $B^\oplus$ and $B^\ominus$ are mutually singular, $b^\oplus_{\overline{\imath}, \overline{\jmath}} = 0$, and  this is absurd, since from Proposition \ref{carperm+-} we have that $b^\oplus_{\overline{\imath}, \overline{\jmath}} = \sum_{h \in X^Y} r(h)_{\overline{\imath}, \overline{\jmath}} \mu^\oplus(h) \ge \mu^\oplus(f)> 0$, because invariance property of $\mu^\oplus$ implies that $\mu^\oplus\left(\overline{h}\right) = \mu^\oplus(f)$.
% \end{proof}

% {\color{orange} Per maggior chiarezza, riscriverei così la dimostrazione della Proposizione  \ref{prop:singular}:
\begin{proof}
For every $h \in X^{Y}$, we denote by $r_{ij}(h)$ the entry in position $(i,j)$ of the matrix $R(h)$. Suppose, for a contradiction, that $\mu^\oplus$ and $\mu^\ominus$ are not mutually singular, i.e., there exists a function $\bar h \in X^{Y}$ such that $\mu^{\oplus}(\bar h) > 0$ and $\mu^{\ominus}(\bar h) > 0$. Then, by the definition of $\mu^{\oplus}$ and the inequality $\mu^{\oplus}(\bar h) > 0$, there exists $f \in \mathcal{O}(\bar h)$ such that $c^{\oplus}(f) > 0$. Since $R(f)$ is a rectangular permutation matrix, for every $i \in \underline{m}$ there exists exactly one index $j_i \in \underline{n}$ such that $r_{i j_i}(f) = 1$. Hence, by Proposition~\ref{prop:sum}, it holds that
\[
b^{\oplus}_{i j_i}
= \sum_{h \in X^{Y}} c^{\oplus}(h)\, r_{i j_i}(h)
\ge c^{\oplus}(f)\, r_{i j_i}(f)
= c^{\oplus}(f) > 0.
\]
We recall that the positivity of $b^\oplus_{ij}$ implies that $b^\ominus_{ij} = 0$ for any $(i, j)\in \underline{m}\times\underline{n}$, by definition of $B^\oplus$ and $B^\ominus$. Therefore, $b^{\ominus}_{i j_i} = 0$.
On the other hand, by Proposition~\ref{carperm+-} we have
\[
b^{\ominus}_{i j_i}
= \sum_{h \in X^{Y}} r_{i j_i}(h)\ \mu^{\ominus}(h)
\ge r_{i j_i}(f)\ \mu^{\ominus}(f)
= \mu^{\ominus}(f).
\]
Since $\mu^\ominus$ is a generalized $T$-permutant measure (Proposition~\ref{prop:gpm}), it follows that $\mu^{\ominus}(f) = \mu^{\ominus}(\bar h)$. As a consequence, $b^{\ominus}_{i j_i} \ge \mu^{\ominus}(\bar h) > 0$, which contradicts the equality $b^{\ominus}_{i j_i} = 0$.
\end{proof}

Proposition~\ref{prop:singular} ensures that $\mu^\oplus$ and $\mu^\ominus$ are the Hahn-Jordan decomposition of $\mu$.
% \subsection{Generalization to the non-homogeneous case}
% Let us consider a group equivariant operator $(F,T)\colon (\R^X,G) \to (\R^Y,K)$. Unlike the previous sections, we have no requirements on the action of $T(G)$ on $Y$.
% Additionally, we assume that
% the cardinality of the orbit space $Y/T(G) = \{Y_\ell\}_{\ell=1}^{\overline{d}}$ is equal to $\overline{d}$. 
% Considering the projection $\pi_\ell\colon \mathbb{R}^Y \to \mathbb{R}^{Y_\ell}$ induced by the inclusion $\imath_\ell\colon Y_\ell \hookrightarrow Y$, we can define the function $F^\ell = \pi_\ell \circ F\colon \mathbb{R}^X \to \mathbb{R}^{Y_\ell}$ for any $\ell \in [\overline{d}]$. Moreover, each bijection $k \in K$ can be appropriately restricted (and corestricted) to a bijection on the orbit $Y_\ell$, and this procedure induces a group homomorphism $\delta_\ell$ that associates each $k \in K$ with its restriction (i.e. $\delta_\ell(k)$ is the unique function such that $ \imath_\ell \circ \delta_\ell(k) =  k \circ \imath_\ell$). Hence, we can define the group homomorphism $T_\ell = \delta_\ell \circ T$. In particular, we obtain that $F = (F_1, \dots, F_{\overline{d}})$ and $T = ( T_1, \dots, T_{\overline{d}})$.
% It can be easily seen that $(F_\ell,  T_\ell)$ satisfies the hypotheses of Theorem~\ref{teo:linear}, and therefore for every $\ell \in [\overline{d}]$ there exists a $T_\ell$-permutant measure $\mu_\ell$ on $X^{Y_\ell}$ such that $F_\ell(\varphi) = F_{\mu_\ell}(\varphi) = \sum_{h_\ell \in X^{Y_\ell}} \varphi h_\ell \mu_\ell(h_\ell)$ for every $\varphi \in \mathbb{R}^X$.

\section{Linear GENEOs and $T$-permutant measures}
\label{Linear}
\sectionmark{Building linear GENEOs}

We now extend Theorem~\ref{maint}
to include non-expansivity. To this end, one
final step is required, in the form of Lemma~\ref{lemma:ne}. This result
allows us to connect the generalized $T$-permutant measure $\mu$ with the norms of $F$
and $\p$ under the same hypotheses as in Theorem~\ref{maint}.


% {\color{lightgray}\begin{proposition}
%     \label{lemma:ne}
%     \textcolor{purple}{N: Correggere "surjective" con le nuove ipotesi.} 
%     Assume that $G \subseteq \mathrm{Aut}(X), K \subseteq \mathrm{Aut}(Y)$, $K$ transitively acts on the finite set $Y$, $T \colon G \to K$ is a surjective homomorphism and $F$ is a map from $\mathbb{R}^X$ to $\mathbb{R}^Y$. It holds that 
%     \[
%     \sum_{h \in X^Y} \lvert \mu (h) \rvert = \max_{\varphi \in \R^X \setminus \{\mathbf{0}\}}\frac{\left\| F(\varphi)\right\|_\infty}{\left\| \varphi\right\|_\infty}.
%     \]
% \end{proposition}
% }
\begin{lemma}
    \label{lemma:ne}
    Assume that $G \subseteq \mathrm{Aut}(X), K \subseteq \mathrm{Aut}(Y)$, $T \colon G \to K$ is a homomorphism such that $T(G)$ transitively acts on $Y$, and $F$ is a $T$-equivariant map from $\mathbb{R}^X$ to $\mathbb{R}^Y$.  
    Let $\mu$ be the measure associated with $F$ in the proof of Theorem~\ref{maint}.
    Then, it holds that 
    \[
    \sum_{h \in X^Y} \lvert \mu (h) \rvert = \max_{\varphi \in \R^X \setminus \{\mathbf{0}\}}\frac{\left\| F(\varphi)\right\|_\infty}{\left\| \varphi\right\|_\infty}.
    \]
\end{lemma}

\begin{proof}
The statement is trivially true if $F$ is the zero operator, since in this case
$\mu$ coincides with the zero measure on $X^Y$, because
Proposition~\ref{prop:sum} shows that $c^\oplus(h)$ and $c^\ominus(h)$ vanish for
every $h \in X^Y$. Hence, we may assume that $F$ is not the zero operator and
that $B$ is not the zero matrix.
Setting $c:= c^\oplus - c^\ominus$, Corollary~\ref{corc} implies that $\lvert c(h) \rvert = c^\oplus(h) + c^\ominus(h)$ for every $h \in X^Y$. By definition of $\mu^\oplus$ and $\mu^\ominus$, we have that, for any $h \in X^Y$:
    \begin{align*}
    \sum_{f \in \mathcal{O}(h)} \mu^\oplus(f) & = \sum_{f \in \mathcal{O}(h)} c^\oplus(f)
    , \\
    \sum_{f \in \mathcal{O}(h)} \mu^\ominus(f) & = \sum_{f \in \mathcal{O}(h)} c^\ominus(f).
    \end{align*}
    It follows that, for each $h \in X^Y$,
    \begin{align*}
    \sum_{f \in \mathcal{O}(h)} \left| \mu(f) \right| & = \sum_{f \in \mathcal{O}(h)} \mu^\oplus(f) + \sum_{f \in \mathcal{O}(h)} \mu^\ominus(f) \\
    & = \sum_{f \in \mathcal{O}(h)} c^\oplus(f) + \sum_{f \in \mathcal{O}(h)} c^\ominus(f) \\
    & = \sum_{f \in \mathcal{O}(h)} \left| c(f) \right|,
    \end{align*}
    and hence,
    \begin{equation}
    \label{eq:negeq}
    \sum_{h \in X^Y} \left| c(h) \right| = \sum_{h \in X^Y} \left| \mu(h) \right|.
    \end{equation}
    Let us set $\mathds{1}_X := \sum_{r=1}^n \mathds{1}_{x_r}$ and $\mathds{1}_Y := \sum_{s=1}^m \mathds{1}_{y_s}$. 
    By observing that 
    $R(h) [\mathds{1}_X]=[\mathds{1}_Y]$ for every $h\in X^Y$, from Proposition~\ref{prop:sum} we obtain that 
    \begin{align*}
[F^\oplus\left(\mathds{1}_X\right)]&=\sum_{h \in X^Y} c^\oplus(h) R(h) [\mathds{1}_X] = \left(\sum_{h \in X^Y}c^\oplus(h)\right) [\mathds{1}_Y],
    \\ [F^\ominus\left(\mathds{1}_X\right)]&=\sum_{h \in X^Y} c^\ominus(h) R(h) [\mathds{1}_X] = \left(\sum_{h \in X^Y}c^\ominus(h)\right) [\mathds{1}_Y].
    \end{align*}
    Since $F^\oplus$ (resp. $F^\ominus$) is a $T$-equivariant linear map, any row of $B^\oplus$ (resp. $B^\ominus$) is a permutation of the first row (Lemma~\ref{lempermutationrow}). Then, we get
    \begin{align*}
    [F^\oplus(\mathds{1}_X)] & = B^\oplus [\mathds{1}_X] = \left( \sum_{j=1}^n b^\oplus_{1j}\right) [\mathds{1}_Y]
    , \\
    [F^\ominus(\mathds{1}_X)] & = B^\ominus [\mathds{1}_X] = \left( \sum_{j=1}^n b^\ominus_{1j}\right) [\mathds{1}_Y].
    \end{align*}
    It follows that
    \begin{align*}
    \sum_{h\in X^Y}c^\oplus(h) & = \sum_{j=1}^n b^\oplus_{1j}
    , \\
    \sum_{h\in X^Y}c^\ominus(h) & = \sum_{j=1}^n b^\ominus_{1j}.
    \end{align*}
Therefore,
    \begin{align}
    \label{eq:neeq}
    \sum_{h \in X^Y} \left| c(h) \right| & = \sum_{h \in X^Y} c^\oplus(h) + \sum _{h \in X^Y} c^\ominus(h) \nonumber \\
    & = \sum_{j=1}^n b^\oplus_{1j} + \sum_{j=1}^n b^\ominus_{1j} \nonumber \\
    & = \sum_{j=1}^n \left| b_{1j} \right|.
    \end{align}
{\color{black} Moreover, recalling once again that any row of $B$ is a permutation of the first
row of $B$,} we obtain that 
%\todo{F: io qui metterei un $\max_i$ nella disuguaglianza e sostituirei con un $=$, altrimenti non è chiaro sotto sul perché c'è l'uguaglianza.}
    \[
    \sum_{j=1}^n \left| b_{1j} \right| = \left| \sum_{j=1}^n b_{1j} \mathrm{sgn}(b_{1j}) \right| = \max_{i\in\underline{m}} \left| \sum_{j=1}^n b_{ij} \mathrm{sgn}(b_{1j}) \right|.
    \]
    Let us now consider $\overline{\varphi} = \sum^n_{j=1}\mathrm{sgn}(b_{1j})\mathds{1}_{x_j} \in \R^X$. Since \(B\) is not the zero matrix, it follows that $\bar \p \neq\mathbf{0}$ and $\lVert \overline{\varphi} \rVert_\infty = 1$. Moreover, we have
    \begin{align*}
        [F (\overline{\varphi})] & = B [\overline{\varphi}] \\
        & = B \sum^n_{j=1}\mathrm{sgn}(b_{1j})[\mathds{1}_{x_j}] \\
        & = \sum^n_{j=1}\mathrm{sgn}(b_{1j})B[\mathds{1}_{x_j}] \\
        & = \sum^n_{j=1}\mathrm{sgn}(b_{1j})\sum^m_{i=1}b_{ij}[\mathds{1}_{y_i}] \\
        & = \sum^m_{i=1}\left(\sum^n_{j=1}\mathrm{sgn}(b_{1j})b_{ij}\right)[\mathds{1}_{y_i}].
    \end{align*} 
    Hence, from Equations (\ref{eq:negeq}) and (\ref{eq:neeq}) we have 
    \begin{align*}
    \left\| F (\overline{\varphi}) \right\|_\infty & = \left\| \sum^m_{i=1}\left(\sum^n_{j=1}\mathrm{sgn}(b_{1j})b_{ij}\right)\mathds{1}_{y_i} \right\|_\infty \\
    & = \left\|\left(\sum^n_{j=1}\left| b_{1j} \right| \right)\mathds{1}_{Y}\right\|_\infty\\
    & = \sum^n_{j=1}\left| b_{1j} \right| \\
    & = \sum_{h \in X^Y} \left| c(h) \right| = \sum_{h \in X^Y} \left| \mu(h) \right|.
    \end{align*}
It follows that 
$\frac{ \left\| F (\overline{\varphi}) \right\|_\infty}{\left\| \overline{\varphi} \right\|_\infty}= \left\| F (\overline{\varphi}) \right\|_\infty = \sum_{h \in X^Y} \left| \mu(h) \right|$. In particular,
    \[
    \max_{\varphi\in\R^X\setminus\{\mathbf{0}\}}\frac{\left\| F(\varphi)\right\|_\infty}{\left\|\varphi\right\|_\infty} \ge \frac{ \left\| F (\overline{\varphi}) \right\|_\infty}{\left\| \overline{\varphi} \right\|_\infty}= \sum_{h \in X^Y} \left| \mu(h) \right|.
    \]
    Moreover, from Theorem~\ref{maint} we have that, for every $\varphi\in\R^X, F(\varphi)=\sum_{h\in X^Y}\varphi h\ {\mu}(h)$. Hence,
    \begin{align*}
        \left\| F(\varphi) \right\|_\infty & \le \sum_{h\in X^Y}\left\|\varphi h\right\|_\infty \left|{\mu}(h)\right| \\
        & \le \left\|\varphi\right\|_\infty\sum_{h\in X^Y} \left|{\mu}(h)\right|.
    \end{align*}
    Therefore, $\frac{\left\| F(\varphi)\right\|_\infty}{\left\|\varphi\right\|_\infty}\le\sum_{h\in X^Y}{\left|\mu(h)\right|}$ for every $\varphi\in\R^X\setminus\{\mathbf{0}\}$. In particular, the same applies for the maximum over $\p \in \R^X \setminus\{\mathbf{0}\}$. In conclusion, 
    \[
    \sum_{h\in X^Y}{|\mu(h)|}=\max_{\varphi\in\R^X\setminus\{\mathbf{0}\}}\frac{\left\| F(\varphi)\right\|_\infty}{\left\|\varphi\right\|_\infty}.
    \]
\end{proof}

Finally, we are able to state and prove the main result of this section, namely the representation theorem for linear GENEOs.

% \begin{theorem}
%     \label{teo:linear}
%     Assume that $G \subseteq \mathrm{Aut}(X), K \subseteq \mathrm{Aut}(Y)$, $K$ transitively acts on the finite set $Y$, $T \colon G \to K$ is a surjective homomorphism and $F$ is a map from $\mathbb{R}^X$ to $\mathbb{R}^Y$. The map $F$ is a linear group equivariant non-expansive operator from $(\mathbb{R}^X, G)$ to $(\mathbb{R}^Y, K)$ with respect to the homomorphism $T$ if and only if a generalized permutant measure $\mu$ exists such that $F(\varphi) = \sum_{h \in X^Y} \varphi h \mu(h)$ for every $\varphi \in \mathbb{R}^X$ and $\sum_{h \in X^Y} \left| \mu(h) \right| \le 1$.
% \end{theorem}

\begin{theorem}
    \label{teo:linear}
    Assume that $G \subseteq \mathrm{Aut}(X), K \subseteq \mathrm{Aut}(Y)$, $T \colon G \to K$ is a group homomorphism, $T(G) \subseteq K$ transitively acts on the finite set $Y$ and $F$ is a map from $\mathbb{R}^X$ to $\mathbb{R}^Y$. The map $(F,T)$ is a linear GENEO from $(\mathbb{R}^X, G)$ to $(\mathbb{R}^Y, K)$ if and only if a generalized $T$-permutant measure $\mu$ exists such that $F(\varphi) = \sum_{h \in X^Y} \varphi h \ \mu(h)$ for every $\varphi \in \mathbb{R}^X$ and $\sum_{h \in X^Y} \left| \mu(h) \right| \le 1$.
\end{theorem}

\begin{proof}
    Let us assume that $(F,T)$ is a GENEO from $\left( \R^X, G \right)$ to $\left( \R^Y, K \right)$ with respect to $T$. Then, Theorem~\ref{maint} guarantees that a generalized $T$-permutant measure $\mu$ exists such that $F(\p) = \sum_{h \in X^Y} \p h \mu(h)$ for every $\p \in \R^X$. Moreover, Lemma~\ref{lemma:ne} guarantees that $\sum_{h \in X^Y} \left| \mu(h) \right| = \max_{\varphi \in \R^X \setminus \{\mathbf{0}\}} \frac{\left\| F(\p) \right\|_\infty}{\left\| \p \right\|_\infty}$. Since $F$ is non-expansive, $\sum_{h \in X^Y} \left| \mu(h) \right| \le 1$, which is the first implication of the statement. Let us now assume that a generalized $T$-permutant measure $\mu$ exists such that $F(\p) = \sum_{h \in X^Y} \p h \mu(h)$ for every $\p \in \R^X$, with $\sum_{h \in X^Y} \left| \mu(h) \right| \le 1$. Then, Proposition~\ref{prop:GEOperm} guarantees that $F$ is linear and $T$-equivariant. Moreover, we can prove the non-expansivity of $F$:
    \begin{align*}
        \left\| F(\p) \right\|_\infty & = \left\| \sum_{h \in X^Y} \p h \ \mu(h) \right\|_\infty \\
        & \le \sum_{h \in X^Y} \left\| \p h \right\|_\infty \left| \mu(h) \right| \\
        & \le \left\| \p \right\|_\infty \left( \sum_{h \in X^Y} \left| \mu(h) \right| \right) \\
        & \le \left\| \p \right\|_\infty.
    \end{align*}
\end{proof}

Theorem~\ref{teo:linear} represents every linear GENEO between different
perception pairs by means of generalized $T$-permutant measures. As we will show in
the next section, generalized permutant measures are easier to define than
GENEOs; hence, this theorem provides a valuable tool for the construction of
GENEOs. The ability to populate the space of GENEOs with operators that can be
explicitly constructed is fundamental for obtaining a wide range of
applications. For this reason, Theorem~\ref{teo:linear} is not only important
from a theoretical point of view, but also crucial for the applicability of
GENEO theory.


\begin{rem}
    \label{rem:noinj}
    The generalized $T$-permutant measure of Theorem~\ref{teo:linear} is not unique. As an example, let us consider $X = Y = \left\{ 1, 2, 3 \right\}, G = K = \Aut(X)$ and the identity homomorphism $T = \mathrm{id}_{\Aut(X)}$. Let $F \colon \R^X \to \R^X$ be the linear application that maps $\mathds{1}_j$ to $\frac{1}{3}\sum_{i \in X} \mathds{1}_i$ for any $j \in X$. It is easy to check that $F$ is a GENEO with associated matrix
    \[ B = 
    \begin{pmatrix}
        1/3 & 1/3 & 1/3 \\
        1/3 & 1/3 & 1/3 \\
        1/3 & 1/3 & 1/3
    \end{pmatrix}.
    \]
    One could represent $B$ as:
    \[ B = 
    \begin{pmatrix}
        1/3 & 0 & 0 \\
        0 & 1/3 & 0 \\
        0 & 0 & 1/3
    \end{pmatrix} + \begin{pmatrix}
        0 & 1/3 & 0 \\
        0 & 0 & 1/3 \\
        1/3 & 0 & 0
    \end{pmatrix} +     \begin{pmatrix}
        0 & 0 & 1/3 \\
        1/3 & 0 & 0 \\
        0 & 1/3 & 0
    \end{pmatrix} =: \frac{1}{3}R_1 + \frac{1}{3}R_2 + \frac{1}{3}R_3,
    \] and
    \[ B = 
    \begin{pmatrix}
        0 & 0 & 1/3 \\
        0 & 1/3 & 0 \\
        1/3 & 0 & 0
    \end{pmatrix} + \begin{pmatrix}
        0 & 1/3 & 0 \\
        1/3 & 0 & 0 \\
        0 & 0 & 1/3
    \end{pmatrix} +     \begin{pmatrix}
        1/3 & 0 & 0 \\
        0 & 0 & 1/3 \\
        0 & 1/3 & 0
    \end{pmatrix} =: \frac{1}{3}R_4 + \frac{1}{3}R_5 + \frac{1}{3}R_6.
    \]
    One could easily check that $\left\{R_1 \right\}, \left\{R_2, R_3 \right\}$ and $\left\{R_4, R_5, R_6 \right\}$ are three orbits of $X^X$ under the action of $\alpha_T$. We can define these two different generalized $T$-permutant measures:
    \[
    \mu(h) := \begin{cases}
        1/3 & \text{if } P(h) = R_1, R_2, R_3 \\
        0 & \text{otherwise}
    \end{cases},
    \]
    \[
    \nu(h) := \begin{cases}
        1/3 & \text{if } P(h) = R_4, R_5, R_6 \\
        0 & \text{otherwise}
    \end{cases}.
    \]
    Both generalized permutant measures define the same linear GENEO $(F,\mathrm{id}_{\Aut(X)})$.
\end{rem}

\section{Compactness and convexity of the space of linear GENEOs}
\label{Compactness}
% This section aims to establish several results concerning the topology and geometry of the space of linear GENEOs, in particular compactness and convexity.
% \textcolor{red}{F: Io parlerei in tutta la sezione di politopo, che è una definizione più stringente di convesso.} 
% Given a fixed homomorphism $T:G\to K$, we denote with $\mathcal{F}_\mathrm{lin}^T$ the space of all linear GENEOs $\left(F, T \right) \colon \left(\R^X, G\right) \to \left(\R^Y, K\right)$. We highlight the fact that, since $X^Y$ is finite, there is a finite number of orbits of the action $\alpha_T$ which are denoted with $\mathcal{O}_1, \dots, \mathcal{O}_k \subseteq X^Y$. For each $i=1, \dots, k$ we define the signed measures
% \[
%     \nu^+_i(h) := \begin{cases}
%         1/\left| \mathcal{O}_i \right| & \text{if } h \in \mathcal{O}_i \\
%         0 & \text{otherwise}
%     \end{cases}, \quad
%     \nu^-_i(h) := \begin{cases}
%         -1/\left| \mathcal{O}_i \right| & \text{if } h \in \mathcal{O}_i \\
%         0 & \text{otherwise}
%     \end{cases}.
% \]
% We observe that, for every $i$, $\nu^+_i$ and $\nu^-_i$ are generalized permutant measures, since they are constant on the orbits. Moreover, $\sum_{h\in X^Y} \left| \nu_i^+ (h) \right| = \sum_{h\in X^Y} \left| \nu_i^- (h) \right| = 1$. Before continuing, we are going to need the following lemma.

% \begin{lemma}
% \label{lem:convex}
%     Given a set $S = \left\{ x_1, \dots, x_n, -x_1, \dots, -x_n \right\} \subseteq \R^m$ symmetric with respect to the origin, we consider the convex hull $\mathrm{Conv}(S)$ and the set $\mathcal{C}$ of all points of the form $\sum_{i=1}^n a_i x_i + \sum_{i=1}^n b_i (-x_i)$, where $a_i, b_i \ge 0$ for every $i$ and $\sum_{i=1}^n a_i + \sum_{i=1}^n b_i \le 1$. Then $\mathrm{Conv}(S) = \mathcal{C}$.
% \end{lemma}
% \begin{proof}
% Clearly, $\mathrm{Conv}(S) \subseteq \mathcal{C}$, hence to conclude the proof it suffice to show that $\mathcal{C} \subseteq \mathrm{Conv}(S)$. Now we consider $ \mathcal{C} \ni x = \sum_{i=1}^n a_i x_i + \sum_{i=1}^nb_i(-x_i)$. Let us define $\tilde{a}_i = a_i + \frac{1}{2n}\left( 1 - \sum_{j=1}^na_j - \sum_{j=1}^nb_j \right)$, $\tilde{b}_i = b_i + \frac{1}{2n}\left( 1 - \sum_{j=1}^na_j - \sum_{j=1}^nb_j \right)$ for every $i=1, \dots, n$. Then, it holds that
%     \begin{align*}
%         \sum_{i=1}^n \tilde{a}_i + \sum_{i=1}^n \tilde{b}_i & = a_1 + \frac{1 - \sum\limits_{j=1}^na_j - \sum\limits_{j=1}^nb_j}{2n} + \dots + b_n + \frac{1 - \sum\limits_{j=1}^na_j - \sum\limits_{j=1}^nb_j}{2n} \\
%         & = \sum_{i=1}^n a_i + \sum_{i=1}^n b_i + 1 - \sum_{j=1}^n a_j - \sum_{j=1}^n b_j \\
%         & = 1.
%     \end{align*}
%     Moreover, we have that
%     \begin{align*}
%          &\sum_{i=1}^n \tilde{a}_ix_i + \sum_{i=1}^n \tilde{b}_i(-x_i) = \\ & = \sum_{i=1}^n \left( a_i + \frac{1 - \sum\limits_{j=1}^na_j - \sum\limits_{j=1}^nb_j}{2n} \right) x_i  + \sum_{i=1}^n \left( b_i + \frac{1 - \sum\limits_{j=1}^na_j - \sum\limits_{j=1}^nb_j}{2n} \right) (-x_i) 
%         \\& = \sum\limits_{i=1}^n a_i x_i + \frac{\sum\limits_{i=1}^n x_i}{2n} - \frac{\sum\limits_{i=1}^n\sum\limits_{j=1}^n a_j x_i}{2n} - \frac{\sum\limits_{i=1}^n\sum\limits_{j=1}^n b_j x_i}{2n} 
%         \\&- \sum\limits_{i=1}^n b_i x_i - \frac{\sum\limits_{i=1}^n x_i}{2n} + \frac{\sum\limits_{i=1}^n\sum\limits_{j=1}^n a_j x_i}{2n} + \frac{\sum\limits_{i=1}^n\sum\limits_{j=1}^n b_j x_i}{2n} 
%         \\ & = \sum\limits_{i=1}^n a_i x_i + \sum\limits_{i=1}^n b_i (-x_i) = x.
%     \end{align*}
%     Hence, $x \in \mathrm{Conv}(S)$ that concludes the proof.
% \end{proof}
% Let $\mu$ be a generalized $T$-permutant measure. For every $i = 1, \dots, k$ we set
%     \[
%     \lambda_i^+ = \mu(\mathcal{O}_i), \lambda_i^- = 0 \quad \text{if } \mu(\mathcal{O}_i) \geq 0,
%     \]
%     \[
%     \lambda_i^- = \lvert \mu(\mathcal{O}_i)\rvert, \lambda_i^+ = 0 \quad \text{if } \mu(\mathcal{O}_i) < 0.
%     \]
% Hence, we can write $\mu = \sum_{i=1}^k \lambda_i^+ \nu_i^+ + \lambda_i^-\nu_i^-$. Indeed, assuming that $h \in X^Y$ and $h$ belongs to the orbit $\mathcal{O}_{\overline{\imath}}$ for an index $\overline{\imath} \in \{1, \dots,k\}$, we know that $\mu(h) = \mu(\mathcal{O}_{\overline{\imath}})/\lvert \mathcal{O}_{\overline{\imath}}\rvert$ since $\mu$ is constant on the orbits by the definition of generalized $T$-permutant measure. Then, the definitions of $\lambda_i^+, \lambda_i^-\nu_i^+$, and $\nu_i^-$ imply that $\mu(h) = \mu(\mathcal{O}_{\overline{\imath}})/\lvert \mathcal{O}_{\overline{\imath}}\rvert = \lambda_{\overline{\imath}}^+\nu_{\overline{\imath}}^+(h) + \lambda_{\overline{\imath}}^-\nu_{\overline{\imath}}^-(h) = \sum_{i=1}^k \lambda_i^+ \nu_i^+(h) + \lambda_i^- \nu_i^-(h)$.
% \begin{proposition}
% \label{prop:convex}
%     Each generalized $T$-permutant measure $\mu$ such that $\sum_{h \in X^Y}\left| \mu(h) \right| \le 1$ is a convex combination of $\nu^+_i$ and $\nu^-_i$.
% \end{proposition}
% \begin{proof}
%     % For any $i=1, \dots, k$ let us define $\mu_i := \mu(h)$, where $h \in \mathcal{O}_i$ (it is well defined since $\mu$ is constant on the orbits by definition). Since it holds that $\sum_{h \in X^Y}\left| \mu(h) \right| \le 1$, it must also hold
%     % \begin{equation}
%     % \label{eq:compact}
%     % \sum_{i=1, \dots, k}\left| \mu_i \right|\left|\mathcal{O}_i\right| \le 1.
%     % \end{equation}
%     % For every $i = 1, \dots, k$ let $h_i \in \mathcal{O}_i$. We define $\lambda_i^+, \lambda_i^- \in \R$ to be such that
%     % \[
%     % \mu_i = \lambda_i^+\nu_i^+(h_i), \lambda_i^- = 0 \quad \text{if } \mu_i \geq 0,
%     % \]
%     % \[
%     % \mu_i = \lambda_i^-\nu_i^-(h_i), \lambda_i^+ = 0 \quad \text{if } \mu_i < 0.
%     % \]
%     % It is straightforward to prove that $\nu = \sum_{i=1, \dots, k} \lambda_i^+ \mu_i^+ + \lambda_i^-\mu_i^-$. (\textcolor{magenta}{Per vederlo, preso $h \in X^Y$, sia $\mathcal{O}_i$ la sua orbita. Allora $\nu(h) = \lambda_i^+\mu_i^+(h) + \lambda_i^-\mu_i^-(h) = \sum_{i=1, \dots, k} \lambda_i^+ \mu_i^+(h) + \lambda_i^- \mu_i^-(h)$ perché fuori da $\mathcal{O}_i$ tutte le altre $\mu_j(h) = 0$, mentre la prima uguaglianza segue dalla definizione di $\lambda_i^+, \lambda_i^-$}). 
    
    
%     Since $\lambda_i^+, \lambda_i^- \geq 0$ by definition, to conclude the proof it suffice to show that $\sum_{i=1}^k \lambda_i^+ + \lambda_i^- \le 1$, and the statement immediately follows from Lemma~\ref{lem:convex}. It holds that
%     \begin{align*}
%         \sum_{i=1}^k \lambda_i^+ + \lambda_i^- = \sum_{i=1}^k\left| \mu(\mathcal{O}_i) \right| = \sum_{h \in X^Y}\lvert \mu(h) \rvert \le 1.
%     \end{align*}
% \end{proof}

% % From now on, we are going to denote with $F_i^+$ the GENEO associated to $\nu_i^+$, and similarly for $F_i^-$. 
% % Before proceeding, we recall the following Theorem from \cite{guide2006infinite}. (\textcolor{magenta}{È una generalità eccessiva ma non ho trovato di meglio})
% % \textcolor{purple}{Corollario 5.30}
% % \begin{theorem}
% % \label{teo:convcomp}
% %     In a completely metrizable locally convex space, the closed convex hull of a compact set is compact.
% % \end{theorem}

% %The space $\Gamma$ of generalized permutant measures verifies such theorem (for any metric, since they induce the same topology). 
% %Denoting the power set of $X^Y$ as $\mathcal{P}$,
% We define on the space $\Gamma$ of generalized $T$-permutant measures the topology induced by the norm $\mathrm{TV}(\mu) = \mu^\oplus(X^Y) + \mu^\ominus(X^Y)$.
% %:= \sup_{H \in \mathcal{P}} \left| \mu(H) - \nu(H) \right|$. \textcolor{red}{Francesco: se $\mu(H) = \sum_{h \in H} \mu(\{h\})$, a me non torna come definizione.}We recall that $\Gamma$ is a compact space since it is the convex hull of a finite set. 
% Moreover, we endow the space of linear GENEOs with the usual operator norm defined as

% \[
% \left\| F \right\|_\mathrm{op} := \sup\left\{ \left\| F(\p) \right\|_\infty, \p \in \R^X, \left\| \p \right\|_\infty = 1 \right\}.
% \]
% Before proceeding, we recall that, given a permutant measure $\nu$, we can define an associated GENEO $(F_\nu,T)$ (see Theorem~\ref{teo:linear}).
% \begin{theorem}
%     Given a homomorphism $T \colon G \to K$, the space $\mathcal{F}^T_\mathrm{lin}$ is compact and convex.
%     %the convexhull of  $\{F_{\mu_i^+}, i \in \{1, \dots, k\}\}\cup \{F_{\mu_i^-}, i \in \{1, \dots, k\}\}$.
% \end{theorem}
% \begin{proof}
% Given a homomorphism $T$, let us define the following map: $f \colon \Gamma \to \mathcal{F}^T_\mathrm{lin}, f(\mu) = F_\mu$. Since $\Gamma$ is compact and convex, the fact that $f$ is linear and surjective implies the statement.
% \end{proof}
% % To conclude, it suffice to show that $f$ is continuous. Given $\mu, \nu \in \Gamma$, we have that
% %     \begin{align*}
% %         \left\| f(\mu) - f(\nu) \right\|_\mathrm{op} & = \left\| F_\mu - F_\nu \right\|_\mathrm{op} \\
% %         & = \sup_{\p \in \R^X, \left\| \p \right\|_\infty = 1} \left\| \sum_{h \in X^Y} \p h \mu(h) - \sum_{h \in X^Y} \p h \nu(h) \right\|_\infty \\ 
% %         & = \sup_{\p \in \R^X, \left\| \p \right\|_\infty = 1} \left\| \sum_{h \in X^Y} \left( \mu(h) - \nu(h) \right) \p h \right\|_\infty \\
% %         & \le \sup_{\p \in \R^X, \left\| \p \right\|_\infty = 1} \sum_{h \in X^Y} \left| \mu(h) - \nu(h) \right| \left\| \p h \right\|_\infty \\
% %         & \le \sum_{h \in X^Y} \left| \mu(h) - \nu(h) \right| \\
% %         & \le \sup_{h \in X^Y} \left| \mu(h) - \nu(h) \right| \left|X^Y\right| \\
% %         & = \left|X^Y\right| d_\Gamma(\mu, \nu).
% %     \end{align*}
% %     Hence, $f$ is Lipschitz-continuous with constant $\left| X^Y \right|$ and hence continuous.

% \begin{corollary}
%     The space of linear GENEOs w.r.t. any homomorphism is compact.
% \end{corollary}
% \begin{proof}
%     There is a finite number of homomorphisms, and a finite union of compact sets is compact.
% \end{proof}

% \textcolor{red}{F: congetturo che si potrebbe dire qualcosa di più su questo, ma non so se ne valga la pena. Sicuramente l'operatore nullo $F_0$ è GENEO rispetto a qualsiasi omomorfismo $T$. Per la stessa ragione, gli operatori associati alle misure $\mu_\pm(h) = \pm 1/|X^Y|$ sono GENEO rispetto a qualsiasi omomorfismo $T$. Poiché intersezione di convessi è convesso, il segmento che li collega è GENEO rispetto a qualsiasi omomorfismo, di cui $F_0$ è il punto medio. Per ogni $h \in X^Y$, chiamiamo $m_h$ la cardinalità minima dell'orbita di $h$ al variare dell'omomorfismo $T$. Definiamo le misure $\pm \nu_h$ che valgono $\pm1/m_h$ all'interno dell'orbita di $h$ (orbita rispetto all'omomorfismo che verifica il minimo) e $0$ al di fuori. Queste sono ovviamente misure permutanti rispetto all'omomorfismo che verifica il minimo (che può essere diverso per $h$ diversi). Quello che congetturo è che $\bigcup_{T}\mathcal{F}^T_\mathrm{lin}$ è proprio il politopo convesso con estremi i GENEO associati a $\pm \nu_h$. È utile? Non so.}

% \vskip 0.2cm
% \hrule
% \vskip 0.2cm
% \hrule
% \vskip 0.2cm

% {\color{orange} I risultati descritti in questa sezione mi sembrano rilevanti ma li riscriverei come segue.}

% \vskip 0.2cm
% \hrule
% \vskip 0.2cm
% \hrule
% \vskip 0.2cm

{\color{black} This section aims to establish several results concerning the geometry and
topology of the space of linear GENEOs, in particular compactness and
convexity.

In this section, we fix a homomorphism \(T \colon G \to K\) and let \(\mathcal{F}^T_{\mathrm{lin}}\) denote the space of all linear maps \(F \colon \mathbb{R}^X \to \mathbb{R}^Y\) such that $(F, T) \colon (\R^X, G) \to (\R^Y, K)$ is a GENEO. We denote by \(\Gamma\) the set of all generalized \(T\)-permutant measures and observe that \(\Gamma\) is a real vector space. 
We emphasize that, since $X^Y$ is finite, the action $\alpha_T$
admits only finitely many orbits, which we denote by
$\mathcal{O}_1, \dots, \mathcal{O}_k \subseteq X^Y$.
For each $i \in \underline{k}$, we define the signed measure
\[
\nu_i(h) :=
\begin{cases}
1 & \text{if } h \in \mathcal{O}_i, \\
0 & \text{otherwise}.
\end{cases}
\]
We observe that, for every $i\in \underline{k}$, the measure $\nu_i$ 
is a generalized $T$-permutant measure, since it is constant on
orbits. 
Moreover, the measures $\nu_1,\ldots,\nu_k$ are linearly independent.
%Moreover, $\sum_{h \in X^Y} \lvert \nu_i(h) \rvert=\cancel{|X^Y|}$ \textcolor{blue}{$= |\mathcal{O}_i|$ for every $i=1, \dots, k$}.
We also observe that, for every $\p \in \mathbb{R}^X$,
$\|F_{\nu_i}(\p)\|_\infty
=
\left\|
\sum_{h \in X^Y} \p h \ \nu_i(h)
\right\|_\infty
=
\left\|
\sum_{h \in \mathcal{O}_i} \p h
\right\|_\infty
\le
\sum_{h \in \mathcal{O}_i}
\|\p h\|_\infty
\le
|\mathcal{O}_i|\,\|\p\|_\infty$.

%Before continuing, we are going to need the following lemma.

\begin{proposition}
\label{prop:convexity}
The signed measures $\nu_1, \ldots, \nu_k$ form a basis for $\Gamma$.
% The $T$-permutant measures on $X^Y$ are exactly the signed measures that can be expressed as linear combinations of the signed measures $\nu_1, \ldots, \nu_k$.
    %If $\mu$ is a $T$-permutant measure, then
    %$\mu = \sum_{i=1}^k \mu(\mathcal{O}_i) \nu_i$.
\end{proposition}

\begin{proof}
%The definition of a $T$-permutant measure immediately implies that the set of all $T$-permutant measures is a vector space.
On the one hand, it is clear that each measure of the form $\mu=\sum_{i=1}^k a_i \ \nu_i$
is a $T$-permutant measure on $X^Y$.
On the other hand, if $\mu$ is a generalized $T$-permutant measure on $X^Y$, then for every
index $i\in\underline{k}$ there exists a value $a_i$ such that
$\mu(f)=a_i$ for every $f\in\mathcal O_i$.
Since $\nu_i$ is constant on $\mathcal O_i$, it follows that $\mu = \sum_{i=1}^k a_i \ \nu_i$.
\end{proof}

\begin{corollary}
\label{cor:convexity}
The linear GEOs from $(\R^X,G)$ to $(\R^Y,K)$ are exactly the pairs of maps  
$(\sum_{i=1}^k a_i F_{\nu_i},T)$, with $a_1,\ldots,a_k$ arbitrary real numbers. The linear GENEOs from $(\R^X,G)$ to $(\R^Y,K)$ are exactly the pairs of maps $(\sum_{i=1}^k a_i F_{\nu_i},T)$, with  $\sum_{i=1}^k |a_i||\mathcal{O}_i|\le 1$.
    %If $\mu$ is a $T$-permutant measure, then
    %$\mu = \sum_{i=1}^k \mu(\mathcal{O}_i) \nu_i$.
\end{corollary}

\begin{proof}
Theorem~\ref{maint}  guarantees that the linear GEOs from
$(\mathbb{R}^X, G)$ to $(\mathbb{R}^Y, K)$ are precisely the pairs
$(F_\mu, T)$, where $\mu$ is a generalized $T$-permutant measure on $X^Y$.
The first statement of Corollary~\ref{cor:convexity} follows immediately from
this result and Proposition~\ref{prop:convexity}, by observing that if
$\mu = \sum_{i=1}^k a_i \ \nu_i$, then
$F_\mu = \sum_{i=1}^k a_i F_{\nu_i}$.
%since the orbits $\mathcal{O}_i$ are disjoint.
Moreover, Theorem~\ref{teo:linear} guarantees that the linear GENEOs from
$(\mathbb{R}^X, G)$ to $(\mathbb{R}^Y, K)$ are exactly the pairs
$(F_\mu, T)$ for which $\mu$ is a $T$-permutant measure on $X^Y$
satisfying $\sum_{h \in X^Y} |\mu(h)| \le 1$. The second statement of Corollary~\ref{cor:convexity}
follows from Theorem~\ref{teo:linear} by observing that, if $\mu=\sum_{i=1}^k a_i \ \nu_i$, then 

\[
\begin{aligned}
\sum_{h\in X^Y} |\mu(h)|
&=
\sum_{h\in X^Y} \left|\sum_{j=1}^k a_j \ \nu_j(h)\right|\\
&=\sum_{i=1}^k \sum_{h\in \mathcal{O}_i} \left|\sum_{j=1}^k a_j \ \nu_j(h)\right|\\
&=\sum_{i=1}^k \sum_{h\in \mathcal{O}_i} \left|a_i \right|\\
&=\sum_{i=1}^k |a_i| |\mathcal{O}_i|.
\end{aligned}
\]
%\hrule
%Let us start from the statement about GEOs.
% On the one hand, it is trivial to observe that
% $\left(\sum_{i=1}^k a_i F_{\nu_i}, T\right)$
% is a GEO from $(\mathbb{R}^X, G)$ to $(\mathbb{R}^Y, K)$.
% If $\sum_{i=1}^k |a_i| |\mathcal{O}_i|\le 1$, then
% $\left(\sum_{i=1}^k a_i F_{\nu_i}, T\right)$
% is a GENEO, since
% \[
% \begin{aligned}
% \left\|
% \sum_{i=1}^k a_i F_{\nu_i}(\p_1)
% - \sum_{i=1}^k a_i F_{\nu_i}(\p_2)
% \right\|_\infty
% &=
% \left\|
% \sum_{i=1}^k a_i F_{\nu_i}(\p_1 - \p_2)
% \right\|_\infty \\
% &\le
% \sum_{i=1}^k |a_i|
% \left\|F_{\nu_i}(\p_1 - \p_2)\right\|_\infty \\
% %&= \sum_{i=1}^k |a_i|
% %\left\|(\p_1 - \p_2)h \nu_i(h)\right\|_\infty \\
% &\le
% \sum_{i=1}^k |a_i| |\mathcal{O}_i|\,
% \left\|\p_1 - \p_2\right\|_\infty .
% \end{aligned}
% \]

% On the other hand, if $(F, T)$ is a GEO from $(\mathbb{R}^X, G)$ to $(\mathbb{R}^Y, K)$,
% then Theorem~\ref{teo:linear} implies the existence of a generalized $T$-permutant measure $\mu$
% such that
% $F(\varphi) = \sum_{h \in X^Y} \varphi h \, \mu(h)$
% for every $\varphi \in \mathbb{R}^X$.
% From Proposition~\ref{prop:convexity}, there exists $(a_1, \ldots, a_k) \in \mathbb{R}^k$
% such that
% $\mu = \sum_{i=1}^k a_i \nu_i$,
% and hence
% \begin{align*}
% F(\varphi)
% &= \sum_{h \in X^Y} \varphi h \sum_{i=1}^k a_i \nu_i(h) \\
% &= \sum_{i=1}^k a_i \sum_{h \in X^Y} \varphi h \nu_i(h)
% = \sum_{i=1}^k a_i F_{\nu_i}(\varphi).
% \end{align*}
% If $(F,T)$ is a GENEO, then $F$ is non-expansive and hence, in particular, 
% $\|\sum_{i=1}^k a_i F_{\nu_i}(\varphi)\|_\infty\le \|\p\|_\infty$,
% i.e.,
% $\|\sum_{i=1}^k a_i \sum_{h\in \mathcal{O}_i} \varphi h\|_\infty\le \|\p\|_\infty$ for every $\p\in \R^X$.
% If we set $\p={\mathds{1}}_{x_j}$,
% we obtain the inequality
% $\|\sum_{i=1}^k a_i \sum_{h\in \mathcal{O}_i} {\mathds{1}}_{x_j} h\|_\infty\le \|{\mathds{1}}_{x_j}\|_\infty=1$.






% Moreover,
% $\|F(\varphi)\|_\infty
% \le \sum_{i=1}^k |a_i| \, \|F_{\nu_i}(\varphi)\|_\infty
% \le \sum_{i=1}^k |a_i| \, |\mathcal{O}_i| \, \|\varphi\|_\infty$
% for every $\varphi \in \mathbb{R}^X$.
% Since $F$ is linear, if
% $\sum_{i=1}^k |a_i| \, |\mathcal{O}_i| \le 1$,
% then
% $\|F(\varphi_1 - \varphi_2)\|_\infty
% \le \|\varphi_1 - \varphi_2\|_\infty$
% for every $\varphi_1, \varphi_2 \in \mathbb{R}^X$.
\end{proof}

We define on $\Gamma$ the  
distance $\sum_{h \in X^Y} |\mu_1(h)-\mu_2(h)|$.
%\pat{E' questa la metrica giusta?}
Moreover, we endow $\mathcal{F}_\mathrm{lin}^T$ with the usual operator norm defined as

\[
\left\| F \right\|_\mathrm{op} := \sup\left\{ \left\| F(\p) \right\|_\infty, \p \in \R^X, \left\| \p \right\|_\infty = 1 \right\}.
\]
We recall that, given a generalized $T$-permutant measure $\mu$ such that $\sum_{h\in X^Y}|\mu(h)|\le 1$, 
we can define an associated GENEO $(F_\mu,T)$ by setting 
$F_\mu(\p)=\sum_{h \in X^Y} \varphi h \ \mu(h)$.
%$F(\varphi) = \sum_{h \in X^Y} \varphi h \ \mu(h)$ for every $\varphi \in \mathbb{R}^X$ and $\sum_{h \in X^Y} \left| \mu(h) \right| \le 1$
% (see Proposition~\ref{prop:GEOperm} and Theorem~\ref{teo:linear}).\todo{F: lo toglierei o lo sposterei prima, perché l'abbiamo appena usato nella dimostrazione 6.2.}
\begin{theorem}\label{thmcomp&conv}
    Given a homomorphism $T \colon G \to K$, the space $\mathcal{F}_\mathrm{lin}^T$ is compact and coincides with the convex hull of the set $\left\{\frac{1}{|\mathcal{O}_1|}F_{\nu_1},\ldots, 
    \frac{1}{|\mathcal{O}_k|}F_{\nu_k},
    -\frac{1}{|\mathcal{O}_1|}F_{\nu_1},\ldots, 
    -\frac{1}{|\mathcal{O}_k|}F_{\nu_k}\right\}$.
\end{theorem}
\begin{proof}
Let us define the linear map $f$
that takes each $\mu\in\Gamma$ to the operator $F_\mu$.
% $f \colon \Gamma \to \mathcal{F}^T_{\mathrm{lin}}$\todo{F: qui ci andrebbe lo spazio dei GEO come codominio, che non è definito.
% N: io non introdurrei un'ulteriore notazione, andrei a dirlo a parole.} by setting $f(\mu) = F_\mu$.
%{\it (To be fully rigorous, we need to say a bit more:)}

Since
$$\begin{aligned}
\left\|F_{\mu_1}(\varphi)-F_{\mu_2}(\varphi)\right\|_\infty
&=\left\|\sum_{h\in X^Y} \varphi h\ \mu_1(h)-\sum_{h\in X^Y} \varphi h\ \mu_2(h)\right\|_\infty \\
&=\left\|\sum_{h\in X^Y} \varphi h \ (\mu_1(h)-\mu_2(h))\right\|_\infty \\
&\le \sum_{h\in X^Y} \|\varphi h\|_\infty|\mu_1(h)-\mu_2(h)| \\
&\le \sum_{h\in X^Y} \|\varphi \|_\infty|\mu_1(h)-\mu_2(h)| \\
&= \|\varphi \|_\infty\sum_{h\in X^Y} |\mu_1(h)-\mu_2(h)|, 
\end{aligned}
$$
it follows 
that
$\left\| F_{\mu_1}-F_{\mu_2} \right\|_\mathrm{op}\le \sum_{h\in X^Y} |\mu_1(h)-\mu_2(h)|$.
This implies that the function $f$ is non-expansive, and hence continuous. Let us now set
$\Gamma' := \left\{ \mu \in \Gamma : \sum_{h \in X^Y} |\mu(h)| \le 1 \right\}$, and observe that $\mathcal{F}^T_{\mathrm{lin}} = f(\Gamma')$ (Theorem~\ref{teo:linear}).
Since 
%$\mathcal{F}^T_{\mathrm{lin}} = f(\Gamma')$, 
$\Gamma'$ is compact, and $f$
is continuous, it follows that $\mathcal{F}^T_{\mathrm{lin}}$ is compact as well.


Corollary~\ref{cor:convexity} states that 
the linear GENEOs from $(\R^X,G)$ to $(\R^Y,K)$ are exactly the pairs of maps  
$(\sum_{i=1}^k a_i F_{\nu_i},T)$, with  $\sum_{i=1}^k |a_i|\ |\mathcal{O}_i|\le 1$. We recall that, in a real vector space, $\left\{ \sum_{i=1}^k \lambda_iv_i \mid \sum_{i=1}^k |\lambda_i|\le 1 \right\}=\mathrm{conv}
(v_1,\ldots,v_k,-v_1,\ldots,-v_k)$, where $\mathrm{conv}(S)$ is the convex hull of $S$.
Therefore, the linear GENEOs from $(\R^X,G)$ to $(\R^Y,K)$ are exactly the pairs
$(F,T)$ where $F$ belongs to the convex hull of the set
    $\left\{\frac{1}{|\mathcal{O}_1|}F_{\nu_1},\ldots, 
    \frac{1}{|\mathcal{O}_k|}F_{\nu_k},
    -\frac{1}{|\mathcal{O}_1|}F_{\nu_1},\ldots, 
    -\frac{1}{|\mathcal{O}_k|}F_{\nu_k}\right\}$.
\end{proof}

% {\color{orange}[C'è un'altra possibile dimostrazione. Secondo me vale la pena farlo notare perché il risultato di compattezza per i GENEO lineari usa una distanza diversa da quella usata in passato. Dire che in fondo non abbiamo cambiato nulla è un segno di coerenza.]}

\begin{rem}
It is worth pointing out an alternative proof of the compactness of the space $\mathcal{F}^T_\mathrm{lin}$ in Theorem~\ref{thmcomp&conv}.
Let us consider the compact sets
$\Phi = \{\varphi \in \mathbb{R}^X : \|\varphi\|_\infty = 1\}$ and $\Psi = \{\psi \in \mathbb{R}^Y : \|\psi\|_\infty \le 1\}$.
If $F \in \mathcal{F}^T_{\mathrm{lin}}$, we observe that
$F|_{\Phi}(\Phi) \subseteq \Psi$ (since $F$ is non-expansive). Moreover, we recall the metric
$D_{\mathrm{GENEO}}(F_1|_{\Phi}, F_2|_{\Phi})
:=
\max_{\varphi \in \Phi}
\|F_1|_{\Phi}(\varphi) - F_2|_{\Phi}(\varphi)\|_\infty$ on the set $\mathcal{F}=\{F|_{\Phi}:F\in \mathcal{F}^T_\mathrm{lin}\}$
(cf., e.g., \cite{bergomi2019towards}).
Of course,
$\|F_1 - F_2\|_{\mathrm{op}}
=
D_{\mathrm{GENEO}}(F_1|_{\Phi}, F_2|_{\Phi})$.
This means that the map from $(\mathcal{F}^T_\mathrm{lin},\|\cdot\|_{\mathrm{op}})$
to $(\mathcal{F},D_{\mathrm{GENEO}})$ that sends $F$ to $F|_{\Phi}$
is an isometry.
We already know that the space of GENEOs from
$(\Phi, G)$ to $(\Psi, K)$ is compact with respect to 
$D_{\mathrm{GENEO}}$ (see, e.g., \cite{bergomi2019towards}).
Therefore, the subset $\mathcal{F}$ of this space, being closed, is compact as well. Since 
$\mathcal{F}^T_\mathrm{lin}$ is isometric to 
$\mathcal{F}$, it is compact.
%Moreover, the set $\mathcal{F}^T_{\mathrm{lin}}$ is convex because it is the image of a convex set under the linear map $f$.
\end{rem}

\begin{rem}
    We stress that, in general, the set
\[
\left\{ \tfrac{1}{|\mathcal{O}_1|}F_{\nu_1}, \dots, \tfrac{1}{|\mathcal{O}_k|}F_{\nu_k},
-\tfrac{1}{|\mathcal{O}_1|}F_{\nu_1}, \dots, -\tfrac{1}{|\mathcal{O}_k|}F_{\nu_k} \right\}
\]
does not coincide with the set of vertices of the polytope $\mathcal{F}^T_{\mathrm{lin}}$.
Recall that a point $p$ of a polytope $P$ is called a vertex of $P$ if the equality
$p = \lambda p_1 + (1-\lambda)p_2$,
$\lambda \in ]0,1[,\; p_1,p_2 \in P$,
implies $p = p_1 = p_2$.
Let us now consider the setting of Remark~\ref{rem:noinj}, and the three orbits
\[
\mathcal{O}_1 := \{R_1\}, \qquad
\mathcal{O}_2 := \{R_2, R_3\}, \qquad
\mathcal{O}_3 := \{R_4, R_5, R_6\}
\]
in $X^X$, under the action $\alpha_T$.
Following the construction illustrated at the beginning of this section, we
define the corresponding measures $\nu_1, \nu_2, \nu_3$.
Accordingly, for $i \in \{1,2,3\}$, we have
$F_{\nu_i}(\varphi)
= \sum_{h \in X^Y} \varphi h\, \nu_i(h)
= \sum_{h \in \mathcal{O}_i} \varphi h$.
Since $[F_{\nu_i}(\varphi)]
= \sum_{h \in \mathcal{O}_i} R(h)[\varphi]$,
it is easy to verify that
$[F_{\nu_1}] =
\begin{pmatrix}
1 & 0 & 0 \\
0 & 1 & 0 \\
0 & 0 & 1
\end{pmatrix}$, 
$[F_{\nu_2}] =
\begin{pmatrix}
0 & 1 & 1 \\
1 & 0 & 1 \\
1 & 1 & 0
\end{pmatrix}$, 
and
$[F_{\nu_3}] =
\begin{pmatrix}
1 & 1 & 1 \\
1 & 1 & 1 \\
1 & 1 & 1
\end{pmatrix}$, 
It follows that
$\left[\frac{1}{|\mathcal{O}_1|}F_{\nu_1}\right] =
\begin{pmatrix}
1 & 0 & 0 \\
0 & 1 & 0 \\
0 & 0 & 1
\end{pmatrix}$, 
$\left[\frac{1}{|\mathcal{O}_2|}F_{\nu_2}\right] =
\begin{pmatrix}
0 & \tfrac{1}{2} & \tfrac{1}{2} \\
\tfrac{1}{2} & 0 & \tfrac{1}{2} \\
\tfrac{1}{2} & \tfrac{1}{2} & 0
\end{pmatrix}$, 
and
$\left[\frac{1}{|\mathcal{O}_3|}F_{\nu_3}\right] =
\begin{pmatrix}
\tfrac{1}{3} & \tfrac{1}{3} & \tfrac{1}{3} \\
\tfrac{1}{3} & \tfrac{1}{3} & \tfrac{1}{3} \\
\tfrac{1}{3} & \tfrac{1}{3} & \tfrac{1}{3}
\end{pmatrix}$.
In particular,
$\left[\frac{1}{|\mathcal{O}_3|}F_{\nu_3}\right] = \frac{1}{3}\left[\frac{1}{|\mathcal{O}_1|}F_{\nu_1}\right] + \frac{2}{3}\left[\frac{1}{|\mathcal{O}_2|}F_{\nu_2}\right]$.
Therefore, $F_{\nu_3}$ is not a vertex of the polytope $\mathcal{F}^T_{\mathrm{lin}}$.
\end{rem}
}

% \begin{corollary}
%     The space of linear GENEOs w.r.t. any homomorphism is compact.
% \end{corollary}
% \begin{proof}
%     There is a finite number of homomorphisms, and a finite union of compact sets is compact.
% \end{proof}
% }

% \textcolor{blue}{The set $\left\{ \frac{1}{|\mathcal{O}_1|}F_{\nu_1}, \dots, \frac{1}{|\mathcal{O}_k|}F_{\nu_k}, -\frac{1}{|\mathcal{O}_1|}F_{\nu_1}, \dots, -\frac{1}{|\mathcal{O}_k|}F_{\nu_k} \right\}$ does \textcolor{red}{NOT} coincide with the set of vertices of $\mathcal{F}^T_\mathrm{lin}$. We recall that a vertex $v$ is such that if $v = \lambda x + (1-\lambda)y$, with $\lambda \in (0, 1)$, then $v = x = y$. In the same setting as Remark~\ref{rem:noinj}, it is easy to verify that (with the usual convention GENEO - matrix) $B_{\nu_1} = \begin{pmatrix}
%         1 & 0 & 0 \\
%         0 & 1 & 0 \\
%         0 & 0 & 1
%     \end{pmatrix}$, $B_{\nu_2} = \begin{pmatrix}
%         0 & 1 & 1 \\
%         1 & 0 & 1 \\
%         1 & 1 & 0
%     \end{pmatrix}$ and $B_{\nu_3} = \begin{pmatrix}
%         1 & 1 & 1 \\
%         1 & 1 & 1 \\
%         1 & 1 & 1
%     \end{pmatrix}$. Hence, we have the set of GENEOs $\left\{ \pm F_{\nu_1}, \pm \frac{1}{2}F_{\nu_2}, \pm\frac{1}{3}F_{\nu_3}\right\}$. However, $\frac{1}{3} B_{\nu_3} = \frac{1}{3}B_{\nu_1} + \frac{1}{3}B_{\nu_2} = \frac{1}{3}B_{\nu_1} + \frac{2}{3}\cdot\frac{1}{2}B_{\nu_2}$. That is, $\frac{1}{3}F_{\nu_3} = \frac{1}{3}F_{\nu_1} + \frac{2}{3}(\frac{1}{2}F_{\nu_2})$. This means that there exists two GENEOs ($F_{\nu_1}, \frac{1}{2}F_{\nu_2}$) such that $\frac{1}{3}F_{\nu_3}$ is a convex combination of them and they do not coincide. Hence, $\frac{1}{3}F_{\nu_3}$ is not a vertex.
%}

\section{Application}
\label{Application}
The aim of this section is to show how introducing a layer of linear GENEOs before an autoencoder increases its robustness to noise. This approach builds on prior work that leverages topological information, and specifically persistent homology, within autoencoder architectures \cite{pmlr-v119-moor20a}. Our general framework is motivated by the fact that the computation of persistence diagrams can itself be formalized as a GENEO (see Remark~\ref{rem:PH}). Specifically, we demonstrate that training an autoencoder to reconstruct MNIST digits and then testing it on noisy versions of these digits yields better performance when a layer of linear GENEOs for dimensionality reduction is introduced before the autoencoding architecture. Random toroidal translations are applied to MNIST images to enlarge the dataset, enrich the experimental setting, and exploit the translation equivariance properties of the GENEO operators. To evaluate reconstruction performance, we use a CNN trained to recognize digits on clean MNIST and track how much its performance drops when it receives the noisy MNIST reconstruction provided by the various methods.
% \pat{}
% \todo{Qui c'è un commento vuoto}

\subsection{Mathematical setup}
Let $X = \mathbb{Z}_p \times \mathbb{Z}_p$, where $p$ is prime, and let $Y = \mathbb{Z}_p$. Throughout this section, all operations are performed modulo $p$. Let $G$ be the group of translations of $X$, that is, $G = \{g_v(q) = q + v$ for $v \in X\}$. Let $K$ be the group of translations of $Y$, that is, $K = \{k_c(z) = z + c$ for $c \in Y\}$. We consider the following perception pairs: $(\R^X, G)$ and $(\R^Y, K)$. We define the following bilinear form $\cdot : X \times X \to \mathbb{Z}_p$ given by:
\[
u \cdot v := u_1v_1 + u_2v_2,
\]
where $u = (u_1, u_2)$ and $v = (v_1, v_2)$. It is easy to check that $\cdot$ is a symmetric bilinear form on $X$.

% \begin{proposition}
%     The operation $\cdot$ is a symmetric bilinear form on $X$.
% \end{proposition}

% \begin{proof}
%     Let $u, v, w \in X$ and $\lambda \in \mathbb{Z}_p$.
    
%     \textbf{Linearity in the first argument:}
%     \begin{align*}
%         (u+v) \cdot w &= (u_1 + v_1, u_2 + v_2) \cdot (w_1, w_2) \\
%         &= (u_1+v_1)w_1 + (u_2 + v_2)w_2 \\
%         &= u_1w_1 + v_1w_1 + u_2w_2 + v_2w_2 \\
%         &= (u_1w_1 + u_2w_2) + (v_1w_1 + v_2w_2) \\
%         &= u \cdot w + v \cdot w.
%     \end{align*}
    
%    \textbf{Linearity in the second argument:}
%    \begin{align*}
%        u \cdot (v + w) &= (u_1, u_2) \cdot (v_1 + w_1, v_2 + w_2) \\
%        &= u_1(v_1+w_1) + u_2(v_2 + w_2) \\
%        &= u_1v_1 + u_1w_1 + u_2v_2 + u_2w_2 \\
%        &= (u_1v_1 + u_2v_2) + (u_1w_1 + u_2w_2) \\
%        &= u \cdot v + u \cdot w.
%    \end{align*}
%    
%    \textbf{Scalar multiplication:}
%    \begin{align*}
%        (\lambda u) \cdot v &= (\lambda u_1, \lambda u_2) \cdot (v_1, v_2) \\
%        &= \lambda u_1 v_1 + \lambda u_2 v_2 \\
%        &= \lambda(u_1v_1 + u_2v_2) \\
%        &= \lambda (u \cdot v) = u \cdot (\lambda v).
%    \end{align*}
%    
%    \textbf{Symmetry:}
%    \begin{align*}
%        u \cdot v &= u_1v_1 + u_2v_2 = v_1u_1 + v_2u_2 = v \cdot u.
%    \end{align*}
% \end{proof}

\begin{rem}
    There exist non-trivial vectors $u \in \mathbb{Z}_p \times \mathbb{Z}_p$ (i.e., $u \neq e_i$ for any standard basis vector $e_i$) such that $u \cdot u = 1$. For example, in $\mathbb{Z}_{11}$, the vector $(3, 5)$ satisfies this property, since $(3, 5) \cdot (3, 5) = 9 + 25 = 34 \equiv 1 \pmod{11}$.
\end{rem}

Let us now fix a vector $\bar{w} = (\bar{w}_1, \bar{w}_2) \in \mathbb{Z}_p \times \mathbb{Z}_p$ such that $\bar{w} \cdot \bar{w} = 1$. We will call $\bar w$ a \textit{unit vector}, and define the group homomorphism $T_{\bar w} \colon G \to K$ by:
\[
T_{\bar w}(g_v) := k_{v\cdot \bar w}.
\]
%\begin{proposition}
%    $T_{\bar w}$ is a group homomorphism.
%\end{proposition}
%\begin{proof}
%    Let $u, v \in X$. Then we have:
%    \begin{align*}
%        T_{\bar w}(g_u \circ g_v) &= T_{\bar w}(g_{u+v}) \\
%        &= k_{(u+v)\cdot \bar w} \\
%        &= k_{u\cdot \bar w + v \cdot \bar w} \\
%        &= k_{u\cdot \bar w} \circ k_{v\cdot \bar w}\\
%        &= T_{\bar w}(g_u) \circ T_{\bar w}(g_v).
%    \end{align*}
%\end{proof}
It is easy to check that $T_{\bar w}$ is a group homomorphism. We set $\bar{w}^\perp = (-\bar{w}_2, \bar{w}_1)$. The reader should note that $\bar{w}^\perp \cdot \bar w = 0$. For every $t \in \mathbb{Z}_p$, we define $h_{\bar{w}}^t \in X^Y$,
\begin{equation}
\label{eq:permutant}
    h_{\bar{w}}^t(z) := z\bar{w} + t\bar{w}^\perp.
\end{equation}

\begin{lemma}\label{lemma:multiple}
    For any $v \in X$, the vector $v - (v \cdot \bar{w})\bar{w}$ is a scalar multiple of $\bar{w}^\perp$.
\end{lemma}

\begin{proof}
    Let $u = v - (v \cdot \bar{w})\bar{w}$. 
    First, we verify that $u \cdot \bar{w} = 0$:
    \begin{align*}
        u \cdot \bar{w} &= (v - (v \cdot \bar{w})\bar{w}) \cdot \bar{w} \\
        &= v \cdot \bar{w} - (v \cdot \bar{w})(\bar{w} \cdot \bar{w}) \\
        &= v \cdot \bar{w} - v \cdot \bar{w} = 0.
    \end{align*}
  We can thus observe that the following homogeneous square linear system in the unknowns $w_1$ and $w_2$ admits $\bar w$ as a nonzero solution:
    \[
    \begin{cases}
        u \cdot w = u_1w_1 + u_2w_2 = 0 \\
        \bar{w}^\perp \cdot w = -\bar{w}_2w_1 + \bar{w}_1w_2 = 0.
    \end{cases}
    \]
It follows directly from elementary linear algebra (cf. Theorem~7 in \cite{hoffmann1971linear}) that $u$ and $\bar{w}^\perp$ are linearly dependent.
\end{proof}

\begin{proposition}
\label{prop:permutant}
    The set $H_{\bar w} = \left\{ h_{\bar w}^t, t \in \Z_p\right\}$ is invariant with respect to the action $\alpha_{T_{\bar{w}}}$.
\end{proposition}
\begin{proof}
    For every $t\in\Z_p$ and $v\in \Z_p\times \Z_p$, we have
    \begin{align*}
        g_v \circ h_{\bar w}^t \circ T_{\bar w}(g_v^{-1})(z) &= g_v \circ h_{\bar w}^t \circ T_{\bar w}(g_{-v})(z) \\
        &= g_v\left(h_{\bar w}^t\left(k_{-v\cdot \bar w}(z)\right)\right) \\
        &= g_v\left(h_{\bar w}^t(z - v \cdot \bar w) \right) \\
        &= g_v \left( (z-v\cdot \bar w)\bar w + t \bar{w}^\perp \right) \\
        &= (z - v\cdot \bar w)\bar w + t \bar{w}^\perp + v \\
        &= z\bar w +t\bar{w}^\perp + (v - \left(v\cdot \bar w)\bar w\right) \\
        &= z\bar w + t\bar{w}^\perp + t^\prime \bar{w}^\perp \\
        &= h_{\bar w}^{t + t^\prime}(z)
    \end{align*}
  for a suitable real number $t'$, because of Lemma~\ref{lemma:multiple}.
Hence, $\alpha_{T_{\bar{w}}}(h_{\bar{w}}^{t}) \in H_{\bar w}$, which concludes the proof.
\end{proof}

Proposition~\ref{prop:permutant} enables the direct definition of a linear GENEO using Theorem~\ref{teo:linear}. Denoting with $\mu_{\bar w}$ the measure $\mu_{\bar w}(h) = 1/|H_{\bar w}|$ if $h \in H_{\bar w}$, $\mu_{\bar w}(h) = 0$ otherwise, it is straightforward to check that $\mu_{\bar w}$ is a generalized $T_{\bar{w}}$-permutant measure. Hence, the operator $F_{\mu_{\bar w}}$ that maps 
$\varphi \mapsto {\sum_{h \in X^Y} \p h \ \mu_{\bar w}(h)}=\sum_{h \in H_{\bar w}} \frac{\varphi h}{|H_{\bar w}|}$ is a linear GENEO.

\subsection{Experimental setup}
To enable the application of GENEO operators, all methods use a zero-padded version of MNIST images resized to 29×29 pixels, where 29 is chosen as the closest prime number to the original dimension. The CNN classifier consists of three convolutional layers (respectively, 32, 64, and 128 filters) with ReLU activations and max pooling, followed by three fully connected layers (respectively, 256, 128, and 10 units) with dropout regularization (0.25). The network was trained on translated MNIST images ($29\times29$ pixels) for $10$ epochs using Adam optimizer with learning rate 0.001 and cross-entropy loss. On clean test data, the classifier achieved an accuracy above 98\%.

In this study, we compare four approaches: three baseline methods from standard literature (autoencoder, convolutional autoencoder, and variational autoencoder) and a GENEO-enhanced version of the convolutional autoencoder that consists in a preliminary GENEO layer. The standard autoencoder (AE, \cite{hinton2006reducing}) consists of an encoder ($841\to256\to128\to32$) and symmetric decoder with ReLU activations. The convolutional autoencoder (CAE, \cite{masci2011stacked}) uses two convolutional layers with max pooling in the encoder, a 64-dimensional bottleneck, and transposed convolutions in the decoder. The variational autoencoder (VAE, \cite{kingma2013auto}) follows the standard architecture with encoder layers ($841\to256\to128$) mapping to 32-dimensional mean and log-variance vectors, using the reparameterization trick 
% \pat{Cos'è il reparameterization trick?}\todo{F: È la differenza tra variational autoencoder e autoencoder classico. Autoencoder classici: Apprendono una rappresentazione compressa (codifica) dei dati in modo deterministico, senza modellare la distribuzione dei dati. Sono ottimi per compressione e denoising.
% Variational Autoencoder (VAE): Apprendono una distribuzione di probabilità (es. Gaussiana) nello spazio latente. Invece di codificare un dato in un singolo punto, lo codificano nei parametri (media e varianza) di una distribuzione, con il reparametrization trick. Questo permette di generare nuovi dati campionando dalla distribuzione latente e decodificando.} 
and trained with reconstruction loss plus KL divergence. For the GENEO-based approach that we refer as GENEO CAE, we generate 28 distinct GENEOs, where 28 is the number of unit vectors in $\Z_{29}\times\Z_{29}$. Each unit vector generates a different linear GENEO using the previously defined strategy. The first GENEO-layer is formed by taking the 28-tuple of GENEO outputs, resulting in a combined $28\times29$ dimensional representation. This GENEO-transformed data serves as input to the convolutional autoencoder architecture. This architecture combines the equivariant feature extraction of GENEOs with the spatial processing capabilities of convolutional networks. All models were trained for 30 epochs with Adam optimizer (learning rate 0.001) and MSE reconstruction loss. To evaluate reconstruction robustness, we corrupt test images with salt \& pepper noise at levels ranging from 0.05 to 0.4, then measure the drop in CNN classification accuracy when the classifier receives reconstructed images from each model instead of the original clean images.

\subsection{Results}

We evaluate the reconstruction performance of the four approaches on MNIST digits corrupted with salt \& pepper noise at levels ranging from 5\% to 40\%.

\begin{figure}[t]
\centering
\includegraphics[width=0.48\textwidth]{Figures/cnn_drop_accuracy.pdf}
\includegraphics[width=0.48\textwidth]{Figures/error_plot.pdf}
\caption{Left: CNN classification accuracy on reconstructed images. Right: Mean Absolute Error (MAE) between reconstructed and clean images as a complementary measure. The GENEO CAE achieves both the highest classification accuracy and, with the sole exception of 40\% noise, the lowest reconstruction error.}
\label{fig:accuracy}
\end{figure}

Figure~\ref{fig:accuracy} shows the CNN classification accuracy (left) and Mean Absolute Error (right) as a complementary reconstruction measure. The GENEO CAE demonstrates superior performance, achieving the highest classification accuracy across all noise levels while also maintaining the lowest MAE in most cases. This dual improvement indicates that the equivariant features extracted by GENEO operators not only preserve semantic content for accurate classification but also enable better pixel-level reconstruction fidelity. At 15\% noise corruption, the GENEO CAE preserves approximately 90\% accuracy while both standard and convolutional autoencoders drop below 80\%, with consistently lower reconstruction error. The complementary MAE plot further confirms the reconstruction quality of the GENEO CAE across the noise spectrum.

%\begin{figure}[t]
%\centering
%\includegraphics[width=0.5\textwidth]{Figures/error_plot.pdf}
%\caption{Mean Absolute Error (MAE) between reconstructed and clean images. While MAE is not the most indicative metric for this task, it provides a complementary measure of reconstruction fidelity.}
%\label{fig:mae}
%\end{figure}

%Figure~\ref{fig:mae} presents the mean absolute error between reconstructed and clean images. Although MAE is not the most indicative metric for assessing reconstruction quality in this context, we include it as a reference measure. The GENEO CAE achieves consistently lower reconstruction error, with the exception of the 40\% noise level.

\begin{figure}[t]
    \centering
    \includegraphics[width=0.48\textwidth]{Figures/reconstructions_image400_saltpepper.pdf}
    \includegraphics[width=0.48\textwidth]{Figures/reconstructions_image500_saltpepper.pdf}
    \caption{Reconstruction results for test images ID 400 (left) and 500 (right) across increasing salt \& pepper noise levels. Columns show original images and reconstructions from AE, CAE, VAE, and GENEO CAE. The GENEO CAE demonstrates superior ability to preserve digit structure under heavy noise corruption.}
    \label{fig:recon_combined}
\end{figure}

Figure~\ref{fig:recon_combined} provides qualitative visualizations of reconstruction performance on two representative test digits. Even at low noise levels, the GENEO approach produces cleaner images, with fewer blurred or artifact-laden reconstructions. As corruption increases, the advantages of the GENEO layer become more evident. Standard methods quickly lose structural integrity, while GENEO continues to produce recognizable digits, although heavily corrupted. These results validate our hypothesis that GENEO operators extract noise-robust equivariant features that enable superior reconstruction under corruption. The consistent improvement across all evaluation metrics—classification accuracy, reconstruction error, and visual quality—demonstrates that leveraging the mathematical structure of the translation group through GENEO preprocessing provides a principled approach to enhancing autoencoder robustness.

\section{Conclusions}
\label{Conclusions}
In this article we have established a new theorem showing that every linear GENEO between two finite perception pairs can be represented as a sum of the input data, reparameterized by maps weighted through a suitable generalized $T$-permutant measure, under the assumption that the group actions are transitive. This result generalizes a previous theorem that required the input and output data spaces to coincide. Such a generalization is not only of theoretical importance but also significantly broadens the applicability of GENEOs in concrete settings. To demonstrate this practical relevance, we concluded the article with an application of the new theorem that enhances the performance of autoencoders.

Future work will address two directions: removing from the theorem the assumption that the domains of the data functions are finite, and developing concrete methods for constructing generalized permutant measures.

% \section{Generalization to the non-homogeneous case}

% Let us consider a GENEO $(F,T)\colon (\R^X,G) \to (\R^Y,K)$. Unlike the previous sections, we have no requirements on the action of $T(G)$ on $Y$.
% Additionally, we assume that
% %that $\mathcal{D} \subseteq Y$ is a fundamental domain (i.e., a subset of the space which contains exactly one point from each of the orbits of an action) for the action of $T(G)$ on $Y$ 
% the cardinality of the orbit space $Y/T(G) = \{Y_\ell\}_{\ell=1}^{\overline{d}}$ is equal to $\overline{d}$. 
% %\textcolor{blue}{Secondo me da qui fino a Considering the projection non serve, e infatti non la stai usando. }
% %In particular, given an  orbit $\mathcal{O}^Y_\ell=\{y_{\ell_1}, \dots y_{\ell_p}\}$ under the action of $T(G)$, we say that $\mathcal{I}_{ \ell}=\{\ell_1, \dots, \ell_p\} \subseteq [n]$ is the index set of $\mathcal{O}^Y_\ell$. Thus, we obtain a partition of the set $[n]$ using the index sets: 
% %\[[n] = \bigcup_{\ell=1}^{\overline{d}} \mathcal{I}_\ell.\]
% %The index sets are pairwise disjoint because orbits are.

% Considering the projection $\pi_\ell\colon \mathbb{R}^Y \to \mathbb{R}^{Y_\ell}$ induced by the inclusion $\imath_\ell\colon Y_\ell \hookrightarrow Y$, we can define the function $F^\ell = \pi_\ell \circ F\colon \mathbb{R}^X \to \mathbb{R}^{Y_\ell}$ for any $\ell \in [\overline{d}]$. Moreover, each bijection $k \in K$ can be appropriately restricted (and corestricted) to a bijection on the orbit $Y_\ell$, and this procedure induces a group homomorphism $\delta_\ell$ that associates each $k \in K$ with its restriction (i.e. $\delta_\ell(k)$ is the unique function such that $ \imath_\ell \circ \delta_\ell(k) =  k \circ \imath_\ell$). Hence, we can define the group homomorphism $T^\ell = \delta_\ell \circ T$.
% It can be easily seen that $(F^\ell,  T^\ell)$ satisfies the hypotheses of Theorem~\ref{teo:linear}, and therefore for every $\ell \in [\overline{d}]$ there exists a $T^\ell$-permutant measure $\mu_\ell$ on $X^{Y_\ell}$ such that $F^\ell(\varphi) = F_{\mu_\ell}(\varphi) = \sum_{h_\ell \in X^{Y_\ell}} \varphi h_\ell \mu_\ell(h_\ell)$ for every $\varphi \in \mathbb{R}^X$. 
% In order to obtain functions that have $Y$ as their domain, for any $\ell \in [\overline{d}]$ we consider the operator $e_\ell \colon \mathbb{R}^{Y_\ell} \to \mathbb{R}^Y$ defined as:
% \[
% e_\ell(\psi)(y) := \begin{cases}
%         \psi(y) & \text{if } y \in Y_\ell\\
%         0 & \text{otherwise}.
%     \end{cases}
% \]
% Then, we have that
% \begin{align*}
% F& = \sum_{\ell =1}^{\overline{d}}e_\ell(F^\ell) 
% \\ & = \sum_{\ell=1}^{\overline{d}} e_\ell \left(\sum_{h_\ell \in X^{Y_\ell}} \varphi \circ h_\ell \cdot \mu_\ell(h_\ell)\right) \\
% & = \sum_{\ell = 1}^{\overline{d}} \sum_{h_\ell \in X^{Y_\ell}}
% e_\ell(\varphi \circ h_\ell) \cdot \mu_\ell(h_\ell)
% \\ & = \sum_{h \in X^{Y}} {\color{green}\sum_{\ell = 1}^{\overline{d}} e_\ell(\varphi \circ h \circ \imath_\ell) \cdot \frac{\mu_\ell(h \circ \imath_\ell)}{m^{n- \lvert Y_\ell \rvert}}}.
% %\\ & = \sum_h \varphi h \star \mu(h)
% \end{align*}
% \textcolor{purple}{N: bisogna spiegare il perchè di quel fratto?}
% We note that if $\overline{d}=1$ we obtain the formula in Theorem~\ref{maint} \textcolor{blue}{F: Idem come sopra}.
% Now, we are ready to define a measure on $X^Y$ as follows 
%  $\mu(h) := \sum_{\ell=1}^{\overline{d}} \frac{ \mu_\ell(h \circ \imath_\ell)}{m^{n- \lvert Y_\ell \rvert}}$ for every $h \in X^Y$.
% \beginlemma
%     $\mu$ is a generalized $T$-permutant mesure.
% \endlemma
% \begin{proof}
%     We consider $h, h' \in X^Y$ such that $h' = g h T(g)$ for some $g \in G$. For every $\ell \in [\overline{d}]$, it holds that $h'\circ \imath_\ell = g (h \circ \imath_\ell) T^\ell(g)$, and hence  $\mu_\ell(h'\circ \imath_\ell) = \mu_\ell(h \circ \imath_\ell)$. It follows immediately that $\mu(h)=\mu(h')$.
% \end{proof}
% When the domain $Y$ does not have a single orbit under the action of the group $T(G)$, we can find a $T$-permutant measure on $X^Y$ as in the previous case, but it cannot be applied in the same way as in the previous sections. In a certain sense, we need to multiply, orbit by orbit, each component of the obtained $T$-permutant measure.  
% \textcolor{red}{In conclusion, it is easy to see that, given a $T$-permutant measure $\mu$ on $X^Y$, it is always possible to directly obtain its components $\mu_\ell$ and derive a representation of it.}
% \textcolor{purple}{N: in questa parte rossa vorrei dire che presa una $\mu$ generica posso sempre scriverla come somma di delle $\mu_\ell$}

\section*{Code availability}
The experiments and applications can be reproduced by installing the open-source Python package, available at \url{https://doi.org/10.5281/zenodo.18230049}.

\section*{Data availability}
No datasets were generated during the current study. The dataset analysed is available at
MNIST~\cite{lecun1998mnist}

\section*{Acknowledgments}
F.C., P.F. and N.Q. conducted a portion of their research within the framework of the CNIT WiLab National Laboratory and the WiLab-Huawei Joint Innovation Center. P.F.'s work received partial support from INdAM-GNSAGA, the COST Action CaLISTA, and the HORIZON Research and Innovation Action PANDORA. F.C. was affiliated with the University of Pisa and CNR-ISTI (Pisa) and with Inria, Centre Inria d'Université Côte d'Azur, during the preparation of this work.

\bibliographystyle{plain}
\bibliography{bib}

\end{document}
